% LLNL-MI-422102-DRAFT

\section{The \texorpdfstring{\MPI/}{MPI} Tool Information Interface}
\mpitermtitleindex{tool information interface}
\label{sec:mpit}

\MPI/ implementations often use internal variables to control their
behavior and performance and rely on internal events for their
implementation.  Understanding and manipulating these variables and
tracking these events
can provide a more efficient execution environment or improve performance
for many applications.  This section describes the \MPI/ tool information
interface, which provides a mechanism for \MPI/ implementors to expose
variables, each of which represents a particular property, setting, or
performance measurement from within the \MPI/ implementation, as
well as expose events that can be tracked by tools.  The
interface is split into three parts: the first part provides information
about, and supports the setting of, control variables through which the \MPI/
implementation tunes its configuration.  The second part provides access to
performance variables that can provide insight into internal performance
information of the \MPI/ implementation. The third part enables tools
to query available events within an \MPI/ implementation and
register callbacks for them.

To avoid restrictions on the \MPI/ implementation, the \MPI/ tool
information interface allows the implementation to specify which control variables,
performance variables, and events exist.  Additionally, the user of the \MPI/ tool
information interface can obtain metadata about each available variable or event,
such as its datatype, and a textual description.  The \MPI/ tool
information interface provides the necessary routines to find all variables and events
that exist in a particular \MPI/ implementation; to query their properties;
to retrieve descriptions about their meaning; to access and, if
appropriate, to alter their values; and (in case of events) set callbacks triggered
by them.

Variables, events, and categories across connected \MPI/ processes with equivalent names
are required to have the same meaning (see the definition of ``equivalent''
as related to strings in Section~\ref{sec:mpit:strings}).  Furthermore,
enumerations with equivalent names across connected \MPI/ processes are required
to have the same meaning, but are allowed to comprise different enumeration
items.  Enumeration items that have equivalent names across connected
\MPI/ processes in enumerations with the same meaning must also have the same
meaning.  In order for variables and categories to have the same meaning,
routines in the tools information interface that return details for those
variables and categories have requirements on what parameters must be
identical.  These requirements are specified in their respective sections.

\begin{rationale}
The intent of requiring the same meaning for entities with equivalent names
is to enforce consistency across connected \MPI/ processes.  For example,
variables describing the number of packets sent on different types of
network devices should have different names to reflect their potentially
different meanings.
\end{rationale}

The \MPI/ tool information interface can be used independently from the
\MPI/ communication functionality.  In particular, the routines of this
interface can be called before \MPIisinitialized/ and
after \MPIisfinalized/.  In order to support this behavior cleanly,
the \MPI/ tool information interface uses separate initialization and
finalization routines.  All identifiers used in the \MPI/ tool information
interface have the prefix \code{MPI\_T\_}.

On success, all \MPI/ tool information interface routines return
\error{MPI\_SUCCESS}, otherwise they return an appropriate and unique
return code indicating the reason why the call was not successfully
completed.  Details on return codes can be found in
Section~\ref{sec:mpit:retcodes}.  However, unsuccessful calls to the \MPI/
tool information interface are not fatal and do not impact the execution of
subsequent \MPI/ routines.

Since the \MPI/ tool information interface primarily focuses on tools and
support libraries, \MPI/ implementations are only required to provide C
bindings for functions and constants introduced in this section.  Except
where otherwise noted, all conventions and principles governing the C
bindings of the \MPI/ API also apply to the \MPI/ tool information
interface, which is available by including the \code{mpi.h} header file.
All routines in this interface have local semantics.

\begin{users}
The number and type of control variables, performance variables, and events can vary
between \MPI/ implementations, platforms and different builds of the same
implementation on the same platform as well as between runs.  Hence, any
application relying on a particular variable will not be portable.
Further, there is no guarantee that the number of variables and variable
indices are the same across connected \MPI/ processes.

This interface is primarily intended for performance monitoring tools,
support tools, and libraries controlling the application's environment.
When maximum portability is desired, application programmers should either
avoid using the \MPI/ tool information interface or avoid being dependent
on the existence of a particular control or performance variable or of a
particular event.
\end{users}



\subsection{Verbosity Levels}
\mpitermtitleindex{verbosity level!tool information interface}
\mpitermindex{tool information interface!verbosity level}
\label{sec:mpit:verbose}

The \MPI/ tool information interface provides access to internal
configuration and performance information through a set of control and
performance variables defined by the \MPI/ implementation.  Since some
implementations may export a large number of variables, variables are
classified by a verbosity level that categorizes both their intended
audience (end users, performance tuners or \MPI/ implementors) and a
relative measure of level of detail (basic, detailed or all).  These
verbosity levels are described by a single integer.
Table~\ref{table:tools:mpit:verbosity} lists the constants for all possible
verbosity levels.  The values of the constants are monotonic in the order
listed in the table; i.e., \mpiconst{MPI\_T\_VERBOSITY\_USER\_BASIC}  $<$
\mpiconst{MPI\_T\_VERBOSITY\_USER\_DETAIL}  $<$ \dots $<$
\mpiconst{MPI\_T\_VERBOSITY\_MPIDEV\_ALL}.

% This could be a constlist
\begin{table}[t]
\caption{\MPI/ tool information interface verbosity levels}
\label{table:tools:mpit:verbosity}
\begin{center}
\setlength{\tabcolsep}{3pt}%%ALLOWLATEX%% % to help fit in page width
\begin{tabular}{|l|l|}
\hline
\mpiconstmain{MPI\_T\_VERBOSITY\_USER\_BASIC} &  Basic information of interest to users\\
\mpiconstmain{MPI\_T\_VERBOSITY\_USER\_DETAIL} & Detailed information of interest to users\\
\mpiconstmain{MPI\_T\_VERBOSITY\_USER\_ALL} &  All remaining information of interest to users\\
\hline
\mpiconstmain{MPI\_T\_VERBOSITY\_TUNER\_BASIC} &  Basic information required for tuning\\
\mpiconstmain{MPI\_T\_VERBOSITY\_TUNER\_DETAIL} &  Detailed information required for tuning\\
\mpiconstmain{MPI\_T\_VERBOSITY\_TUNER\_ALL} & All remaining information required for tuning\\
\hline
\mpiconstmain{MPI\_T\_VERBOSITY\_MPIDEV\_BASIC} &  Basic information for \MPI/ implementors\\
\mpiconstmain{MPI\_T\_VERBOSITY\_MPIDEV\_DETAIL} &  Detailed information for \MPI/ implementors\\
\mpiconstmain{MPI\_T\_VERBOSITY\_MPIDEV\_ALL} &  All remaining information for \MPI/ implementors\\
\hline
\end{tabular}
\end{center}
\end{table}


\subsection{Binding \texorpdfstring{\MPI/}{MPI} Tool Information Interface Variables to \texorpdfstring{\MPI/}{MPI} Objects}
\label{sec:mpit:assoc}

Each \MPI/ tool information interface variable provides access to a
particular control setting or performance property of the \MPI/
implementation.  A variable may refer to a specific \MPI/ object such as a
communicator, datatype, or one-sided communication window, or the variable
may refer more generally to the \MPI/ environment of the process.  Except
for the last case, the variable must be bound to exactly one \MPI/ object
before it can be used.  Table~\ref{table:tools:mpit:assoc} lists all \MPI/
object types to which an \MPI/ tool information interface variable can be
bound, together with the matching constant that \MPI/ tool information
interface routines return to identify the object type. It is erroneous to
bind a control variable, performance variable, or event to a handle that
would not be valid to use as an input argument to another \MPI/ call (excluding
calls to the \MPI/ Tool Information Interface) at the same point of
execution.

\begin{table}[ht]
\caption{Constants to identify associations of variables}
\label{table:tools:mpit:assoc}
\begin{center}
\begin{tabular}{|l|l|}
\hline
\textbf{Constant} & \textbf{\MPI/ object}\\
\hline
\mpiconstmain{MPI\_T\_BIND\_NO\_OBJECT} & N/A; applies globally to entire \MPI/ process\\
\mpiconstmain{MPI\_T\_BIND\_MPI\_COMM} & \MPI/ communicator\\
\mpiconstmain{MPI\_T\_BIND\_MPI\_DATATYPE} & \MPI/ datatype\\
\mpiconstmain{MPI\_T\_BIND\_MPI\_ERRHANDLER} & \MPI/ error handler\\
\mpiconstmain{MPI\_T\_BIND\_MPI\_FILE} & \MPI/ file handle\\
\mpiconstmain{MPI\_T\_BIND\_MPI\_GROUP} & \MPI/ group\\
\mpiconstmain{MPI\_T\_BIND\_MPI\_OP} & \MPI/ reduction operator\\
\mpiconstmain{MPI\_T\_BIND\_MPI\_REQUEST} & \MPI/ request\\
\mpiconstmain{MPI\_T\_BIND\_MPI\_WIN} & \MPI/ window\\
\mpiconstmain{MPI\_T\_BIND\_MPI\_MESSAGE} & \MPI/ message object\\
\mpiconstmain{MPI\_T\_BIND\_MPI\_INFO} & \MPI/ info object\\
\mpiconstmain{MPI\_T\_BIND\_MPI\_SESSION} & \MPI/ session\\
\hline
\end{tabular}
\end{center}
\end{table}


\begin{rationale}
Some variables have meanings tied to a specific \MPI/ object.  Examples
include the number of send or receive operations that use a particular
datatype, the number of times a particular error handler has been called,
or the communication protocol and ``eager limit'' used for a particular
communicator.  Creating a new \MPI/ tool information interface variable for
each \MPI/ object would cause the number of variables to grow without
bound, since they cannot be reused to avoid naming conflicts.  By
associating \MPI/ tool information interface variables with a specific
\MPI/ object, the \MPI/ implementation only must specify and maintain a
single variable, which can then be applied to as many \MPI/ objects of the
respective type as created during the program's execution.
\end{rationale}



\subsection{Convention for Returning Strings}
\label{sec:mpit:strings}

Several \MPI/ tool information interface functions return one or more
strings.  These functions have two arguments for each string to be
returned: an \OUT\ parameter that identifies a pointer to the buffer in
which the string will be returned, and an \INOUT\ parameter to pass the
length of the buffer.  The user is responsible for the memory allocation of
the buffer and must pass the size of the buffer ($n$) as the length
argument.  Let $n$ be the length value specified to the function.  On
return, the function writes at most $n-1$ of the string's characters into
the buffer, followed by a null terminator.  If the returned string's length
is greater than or equal to $n$, the string will be truncated to $n-1$
characters.  In this case, the length of the string plus one (for the
terminating null character) is returned in the length argument.  If the
user passes the null pointer as the buffer argument or passes 0 as the
length argument, the function does not return the string and only returns
the length of the string plus one in the length argument.  If the user
passes the null pointer as the length argument, the buffer argument is
ignored and nothing is returned.

\MPI/ implementations behave as if they have an internal character array
that is copied to the output character array supplied by the user.  Such
output strings are only defined to be equivalent if their notional
source-internal character arrays are identical (up to and including the
null terminator), even if the output string is truncated due to a small
input length parameter $n$.



\subsection{Initialization and Finalization}
\label{sec:mpit:init}

The \MPI/ tool information interface requires a separate set of
initialization and finalization routines.

\begin{funcdef}{MPI\_T\_INIT\_THREAD}{(\mbox{required}, \mbox{provided})}
\funcarg{\IN}{required}{desired level of thread support (integer)}
\funcarg{\OUT}{provided}{provided level of thread support (integer)}
\end{funcdef}
\mpibindmain{MPI\_T\_init\_thread}{(\gb{}int~required, int~*provided)}

All programs or tools that use the \MPI/ tool information interface must
initialize the \MPI/ tool information interface in the processes that will
use the interface before calling any other of its routines.  A user can
initialize the \MPI/ tool information interface by calling
\mpifunc{MPI\_T\_INIT\_THREAD}, which can be called multiple times.  In
addition, this routine initializes the thread environment for all routines
in the \mpi/ tool information interface.  Calling this routine when the
\MPI/ tool information interface is already initialized has no effect
beyond increasing the reference count of how often the interface has been
initialized.  The argument \mpiarg{required} is used to specify the desired
level of thread support.  The possible values and their semantics are
identical to the ones that can be used with \mpifunc{MPI\_INIT\_THREAD}
listed in Section~\ref{sec:ei-threads}.  The call returns in
\mpiarg{provided} information about the actual level of thread support that
will be provided by the \MPI/ implementation for calls to \MPI/ tool
information interface routines.  It can be one of the four values listed in
Section~\ref{sec:ei-threads}.

The \MPI/ specification does not require all \MPI/ processes to exist
before \MPIisinitialized/.  If the \MPI/ tool information
interface is used before initialization of \MPI/, the user is
responsible for ensuring that the \MPI/ tool information interface is
initialized on all processes it is used in.  Processes created by the \MPI/
implementation during initialization inherit the status of the \MPI/
tool information interface (whether it is initialized or not as well as all
active sessions and handles) from the process from which they are created.

Processes created at runtime as a result of calls to \MPI/'s dynamic
process management require their own initialization before they can use the
\MPI/ tool information interface.

\begin{users}
If \mpifunc{MPI\_T\_INIT\_THREAD} is called before
\mpifunc{MPI\_INIT\_THREAD}, the requested and provided thread level for
\mpifunc{MPI\_T\_INIT\_THREAD} may influence the behavior and return value
of \mpifunc{MPI\_INIT\_THREAD}.  The same is true for the reverse order.
Likewise, when using the \spm{} (Section~\ref{sec:model_sessions}), the requested
and provided thread level for \mpifunc{MPI\_T\_INIT\_THREAD} may influence
the behavior and return values of \mpifunc{MPI\_SESSION\_INIT} (see Section~\ref{sec:model_sessions}), with the
same being true for the reverse order.
\end{users}

\begin{implementors}
\MPI/ implementations should strive to make as many control or performance
variables available before \MPI/ initialization (instead of adding them
during initialization) to allow tools the most flexibility.  In
particular, control variables should be available before
\MPI/ initialization if their value cannot be changed after
\MPI/ initialization.
\end{implementors}

\begin{funcdefna}{MPI\_T\_FINALIZE}{()}
\end{funcdefna}
\mpibindmain{MPI\_T\_finalize}{(void)}

This routine finalizes the use of the \MPI/ tool information interface and
may be called as often as the corresponding \mpifunc{MPI\_T\_INIT\_THREAD}
routine up to the current point of execution.  Calling it more times
returns a corresponding return code.  As long as the number of calls to
\mpifunc{MPI\_T\_FINALIZE} is smaller than the number of calls to
\mpifunc{MPI\_T\_INIT\_THREAD} up to the current point of execution, the
\MPI/ tool information interface remains initialized and calls to its
routines are permissible.  Further, additional calls to
\mpifunc{MPI\_T\_INIT\_THREAD} after one or more calls to
\mpifunc{MPI\_T\_FINALIZE} are permissible.

Once \mpifunc{MPI\_T\_FINALIZE} is called the same number of times as the
routine \mpifunc{MPI\_T\_INIT\_THREAD} up to the current point of
execution, the \MPI/ tool information interface is no longer initialized.
The user can reinitialize the interface by a subsequent call
to \mpifunc{MPI\_T\_INIT\_THREAD}.

At the end of the program execution, unless \mpifunc{MPI\_ABORT} is called,
an application must have called \mpifunc{MPI\_T\_INIT\_THREAD} and
\mpifunc{MPI\_T\_FINALIZE} an equal number of times.



\subsection{Datatype System}
\label{sec:types}

All variables managed through the \MPI/ tool information interface
represent their values through typed buffers of a given length and type
using an \MPI/ datatype (similar to regular send/receive buffers).  Since
the initialization of the \MPI/ tool information interface is separate from
the initialization of \MPI/, \MPI/ tool information interface routines can
be called before \MPI/ initialization.  Consequently, these routines can
also use \MPI/ datatypes before \MPI/ initialization.  Therefore, within the
context of the \MPI/ tool information interface, it is permissible to use a
subset of \MPI/ datatypes as specified below before \MPI/ initialization.


\begin{table}[ht]
\caption{\MPI/ datatypes that can be used by the \MPI/ tool information interface}
\label{table:tools:mpit:datatypes}
\begin{center}
\begin{tabular}{l}
\type{MPI\_INT} \\
\type{MPI\_INT32\_T}\\
\type{MPI\_INT64\_T}\\
\type{MPI\_UNSIGNED} \\
\type{MPI\_UNSIGNED\_LONG} \\
\type{MPI\_UNSIGNED\_LONG\_LONG} \\
\type{MPI\_UINT32\_T}\\
\type{MPI\_UINT64\_T}\\
\type{MPI\_COUNT}\\
\type{MPI\_CHAR} \\
\type{MPI\_DOUBLE} \\
\end{tabular}
\end{center}
\end{table}

\begin{rationale}
The \MPI/ tool information interface relies mainly on unsigned datatypes
for integer values since most variables are expected to represent counters
or resource sizes.  \type{MPI\_INT} is provided for additional flexibility
and is expected to be used mainly for control variables and enumeration
types (see below).

Providing all basic datatypes, in particular providing all signed and
unsigned variants of integer types, would lead to a larger number of types,
which tools need to interpret.  This would cause unnecessary complexity in
the implementation of tools based on the \MPI/ tool information interface.
\end{rationale}

The \MPI/ tool information interface only relies on a subset of the basic
\MPI/ datatypes and does not use any derived \MPI/ datatypes.
Table~\ref{table:tools:mpit:datatypes} lists all \MPI/ datatypes that can
be returned by the \MPI/ tool information interface to represent its
variables.

The use of the datatype \type{MPI\_CHAR} in the \MPI/ tool information
interface implies a null-terminated character array, i.e., a string in the
C language.  If a variable has type \type{MPI\_CHAR}, the value of the
\mpiarg{count} parameter returned by \mpifunc{MPI\_T\_CVAR\_HANDLE\_ALLOC}
and \mpifunc{MPI\_T\_PVAR\_HANDLE\_ALLOC} must  be large enough to include
any valid value, including its terminating null character.  The contents of
returned \type{MPI\_CHAR} arrays are only defined from index 0 through the
location of the first null character.

\begin{rationale}
The \MPI/ tool information interface requires a significantly simpler type
system than \MPI/ itself.  Therefore, only its required subset must be
present before \MPI/ initialization and \MPI/ implementations do not need
to initialize the complete \MPI/ datatype system.
\end{rationale}

For variables of type \type{MPI\_INT}, an \MPI/ implementation can provide
additional information by associating names with a fixed number of values.
We refer to this information in the following as an enumeration.  In this
case, the respective calls that provide additional metadata for each
control or performance variable, i.e., \mpifunc{MPI\_T\_CVAR\_GET\_INFO}
(Section~\ref{sec:mpit:cvar}), \mpifunc{MPI\_T\_PVAR\_GET\_INFO}
(Section~\ref{sec:mpit:pvar}), and \mpifunc{MPI\_T\_EVENT\_GET\_INFO}
(Section~\ref{sec:mpit:events}), return a handle of type
\mpictypemain{MPI\_T\_enum} that can be passed to the
following functions to extract additional information.  Thus, the \MPI/
implementation can describe variables with a fixed set of values that each
represents a particular state.  Each enumeration type can have $N$
different values, with a fixed $N$ that can be queried using
\mpifunc{MPI\_T\_ENUM\_GET\_INFO}.

\cdeclindex{MPI\_T\_enum}
\begin{funcdef}{MPI\_T\_ENUM\_GET\_INFO}{(\mbox{enumtype}, \mbox{num}, \mbox{name}, \mbox{name\_len})}
\funcarg{\IN}{enumtype}{enumeration to be queried (handle)}
\funcarg{\OUT}{num}{number of discrete values represented by this enumeration (integer)}
\funcarg{\OUT}{name}{buffer to return the string containing the name of the enumeration item (string)}
\funcarg{\INOUT}{name\_len}{length of the string and/or buffer for \mpiarg{name} (integer)}
\end{funcdef}
\mpibindmain{MPI\_T\_enum\_get\_info}{(\gb{}MPI\_T\_enum~enumtype, int~*num, char~*name, int~*name\_len)}

If \mpiarg{enumtype} is a valid enumeration, this routine returns the
number of items represented by this enumeration type as well as its name.
$N$ must be greater than $0$, i.e., the enumeration must represent at least
one value.


\textoutdesc{name}{the name of the enumeration}

The routine is required to return a name of at least length one.  This name
must be unique with respect to all other names for enumerations that the
\MPI/ implementation uses.


Names associated with individual values in each enumeration
\mpiarg{enumtype} can be queried using \mpifunc{MPI\_T\_ENUM\_GET\_ITEM}.


\cdeclindex{MPI\_T\_enum}
\begin{funcdef}{MPI\_T\_ENUM\_GET\_ITEM}{(\mbox{enumtype}, \mbox{index}, \mbox{value}, \mbox{name}, \mbox{name\_len})}
\funcarg{\IN}{enumtype}{enumeration to be queried (handle)}
\funcarg{\IN}{index}{number of the value to be queried in this enumeration (integer)}
\funcarg{\OUT}{value}{variable value (integer)}
\funcarg{\OUT}{name}{buffer to return the string containing the name of the enumeration item (string)}
\funcarg{\INOUT}{name\_len}{length of the string and/or buffer for \mpiarg{name} (integer)}
\end{funcdef}
\mpibindmain{MPI\_T\_enum\_get\_item}{(\gb{}MPI\_T\_enum~enumtype, int~index, int~*value, char~*name, int~*name\_len)}

\textoutdesc{name}{the name of the enumeration item}

If completed successfully, the routine returns the name/value pair that
describes the enumeration at the specified index.  The call is further
required to return a name of at least length one.  This name must be unique
with respect to all other names of items for the same enumeration.


\subsection{Control Variables}
\mpitermtitleindex{control variable!tool information interface}
\mpitermindex{tool information interface!control variable}
\label{sec:mpit:cvar}

The routines described in this section of the \MPI/ tool information
interface specification focus on the ability to list, query, and possibly
set control variables exposed by the \MPI/ implementation.  These variables
can typically be used by the user to fine tune properties and configuration
settings of the \MPI/ implementation.  On many systems, such variables can
be set using environment variables, although other configuration mechanisms
may be available, such as configuration files or central configuration
registries.  A typical example that is available in several existing \MPI/
implementations is the ability to specify an ``eager limit,'' i.e., an
upper bound on the size of messages sent or received using an eager
protocol.


\subsubsection{Control Variable Query Functions}

An \MPI/ implementation exports a set of $N$ control variables through the
\MPI/ tool information interface.  If $N$ is zero, then the \MPI/
implementation does not export any control variables, otherwise the
provided control variables are indexed from $0$ to $N-1$.  This index
number is used in subsequent calls to identify the individual variables.

An \MPI/ implementation is allowed to increase the number of control
variables during the execution of an \MPI/ application when new variables
become available through dynamic loading.  However, \MPI/ implementations
are not allowed to change the index of a control variable or to delete a
variable once it has been added to the set.  When a variable becomes
inactive, e.g., through dynamic unloading, accessing its value should
return a corresponding return code.

\begin{users}
While the \MPI/ tool information interface guarantees that indices or
variable properties do not change during a particular run of an \MPI/
program, it does not provide a similar guarantee between runs.
\end{users}

The following function can be used to query the number of control
variables, $\mpiarg{num\_cvar}$:

\begin{funcdef}{MPI\_T\_CVAR\_GET\_NUM}{(\mbox{num\_cvar})}
\funcarg{\OUT}{num\_cvar}{returns number of control variables (integer)}
\end{funcdef}
\mpibindmain{MPI\_T\_cvar\_get\_num}{(\gb{}int~*num\_cvar)}

The function \mpifunc{MPI\_T\_CVAR\_GET\_INFO} provides access to
additional information for each variable.

\cdeclindex{MPI\_T\_enum}
\begin{funcdef}{MPI\_T\_CVAR\_GET\_INFO}{(\mbox{cvar\_index}, \mbox{name}, \mbox{name\_len}, \mbox{verbosity}, \mbox{datatype}, \mbox{enumtype}, \mbox{desc}, \mbox{desc\_len}, \mbox{bind}, \mbox{scope})}
\funcarg{\IN}{cvar\_index}{index of the control variable to be queried, value between $0$ and $\mpiarg{num\_cvar}-1$ (integer)}
\funcarg{\OUT}{name}{buffer to return the string containing the name of the control variable (string)}
\funcarg{\INOUT}{name\_len}{length of the string and/or buffer for \mpiarg{name} (integer)}
\funcarg{\OUT}{verbosity}{verbosity level of this variable (integer)}
\funcarg{\OUT}{datatype}{\mpi/ datatype of the information stored in the control variable (handle)}
\funcarg{\OUT}{enumtype}{optional descriptor for enumeration information (handle)}
\funcarg{\OUT}{desc}{buffer to return the string containing a description of the control variable (string)}
\funcarg{\INOUT}{desc\_len}{length of the string and/or buffer for \mpiarg{desc} (integer)}
\funcarg{\OUT}{bind}{type of \mpi/ object to which this variable must be bound (integer)}
\funcarg{\OUT}{scope}{scope of when changes to this variable are possible (integer)}
\end{funcdef}
\mpibindmain{MPI\_T\_cvar\_get\_info}{(\gb{}int~cvar\_index, char~*name, int~*name\_len, int~*verbosity, MPI\_Datatype~*datatype, MPI\_T\_enum~*enumtype, char~*desc, int~*desc\_len, int~*bind, int~*scope)}

After a successful call to \mpifunc{MPI\_T\_CVAR\_GET\_INFO} for a
particular variable, subsequent calls to this routine that query
information about the same variable must return the same information.  An
\MPI/ implementation is not allowed to alter any of the returned values.

If any \OUT\ parameter to \mpifunc{MPI\_T\_CVAR\_GET\_INFO} is a
\code{NULL} pointer, the implementation will ignore the parameter and not
return a value for the parameter.


\textoutdesc{name}{the name of the control variable}

If completed successfully, the routine is required to return a name of at
least length one.  The name must be unique with respect to all other names
for control variables used by the \MPI/ implementation.

The argument \mpiarg{verbosity} returns the verbosity level of the variable
(see Section~\ref{sec:mpit:verbose}).

The argument \mpiarg{datatype} returns the \MPI/ datatype that is used to
represent the control variable.

If the variable is of type \type{MPI\_INT}, \MPI/ can optionally specify
an enumeration for the values represented by this variable and return it in
\mpiarg{enumtype}.  In this case, \MPI/ returns an enumeration identifier,
which can then be used to gather more information as described in
Section~\ref{sec:types}.  Otherwise, \mpiarg{enumtype} is set to
\mpiconstmain{MPI\_T\_ENUM\_NULL}.  If the datatype is not \type{MPI\_INT} or the
argument \mpiarg{enumtype} is the null pointer, no enumeration type is
returned.

\textoutdesc{desc}{a description of the control variable}

Returning a description is optional.  If an \mpi/ implementation does not
return a description, the first character for \mpiarg{desc} must be set to
the null character and \mpiarg{desc\_len} must be set to one at the return
of this call.

The parameter \mpiarg{bind} returns the type of the \MPI/ object to which
the variable must be bound or the value \mpiconst{MPI\_T\_BIND\_NO\_OBJECT}
(see Section~\ref{sec:mpit:assoc}).

The scope of a variable determines whether changing a variable's value is
either local to the \MPI/ process or must be done by the user across multiple
connected \MPI/ processes.  The latter is further split into variables that require changes
in a group of \MPI/ processes and those that require collective changes among all
connected \MPI/ processes.  Both cases can require variables on all participating \MPI/ processes either to be set
to consistent (but potentially different) values or to equal values.
The description provided with the variable
must contain an explanation about the requirements and/or restrictions for
setting the particular variable.

On successful return from \mpifunc{MPI\_T\_CVAR\_GET\_INFO}, the argument
\mpiarg{scope} will be set to one of the constants listed in
Table~\ref{table:tools:mpit:scopes}.

If the name of a control variable is equivalent across connected \MPI/ processes,
the following \OUT\  parameters must be identical:
\mpiarg{verbosity}, \mpiarg{datatype}, \mpiarg{enumtype}, \mpiarg{bind},
and \mpiarg{scope}.  The returned description must be equivalent.

\begin{table}[t]
\caption{Scopes for control variables}
\label{table:tools:mpit:scopes}
\begin{center}
\begin{tabular}{|l|l|}
\hline
\textbf{Scope Constant} & \textbf{Description} \\
\hline
\mpiconstmain{MPI\_T\_SCOPE\_CONSTANT} & read-only, value is constant \\
\mpiconstmain{MPI\_T\_SCOPE\_READONLY} & read-only, cannot be written, but can change \\
\mpiconstmain{MPI\_T\_SCOPE\_LOCAL} & may be writeable, writing only affects the calling \\
& \MPI/ process \\
\mpiconstmain{MPI\_T\_SCOPE\_GROUP} & may be writeable, must be set to consistent values\\
& across a group of connected \MPI/ processes \\
\mpiconstmain{MPI\_T\_SCOPE\_GROUP\_EQ} & may be writeable, must be set to the same value\\
& across a group of connected \MPI/ processes \\
\mpiconstmain{MPI\_T\_SCOPE\_ALL} & may be writeable, must be set to consistent values\\
& across all connected \MPI/ processes \\
\mpiconstmain{MPI\_T\_SCOPE\_ALL\_EQ} & may be writeable, must be set to the same value\\
& across all connected \MPI/ processes \\
\hline
\end{tabular}
\end{center}
\end{table} 

\begin{users}
The \mpiarg{scope} of a variable only indicates if a variable might be
changeable; it is not a guarantee that it can be changed at any time.
\end{users}


\begin{funcdef}{MPI\_T\_CVAR\_GET\_INDEX}{(\mbox{name}, \mbox{cvar\_index})}
\funcarg{\IN}{name}{name of the control variable (string)}
\funcarg{\OUT}{cvar\_index}{index of the control variable (integer)}
\end{funcdef}
\mpibindmain{MPI\_T\_cvar\_get\_index}{(\gb{}const~char~*name, int~*cvar\_index)}

\mpifunc{MPI\_T\_CVAR\_GET\_INDEX} is a function for retrieving the index of a control variable given a known variable name.
The \mpiarg{name} parameter is provided by the caller, and \mpiarg{cvar\_index} is returned by the \MPI/ implementation.
The \mpiarg{name} parameter is a string terminated with a null character.

This routine returns \error{MPI\_SUCCESS} on success and returns
\error{MPI\_T\_ERR\_INVALID\_NAME} if \mpiarg{name} does not match the name
of any control variable provided by the implementation at the time of the
call.

\begin{rationale}
This routine is provided to enable fast retrieval of control variables by a
tool, assuming it knows the name of the variable for which it is looking.
The number of variables exposed by the implementation can change over time,
so it is not possible for the tool to simply iterate over the list of
variables once at initialization.  Although using \MPI/ implementation
specific variable names is not portable across \MPI/ implementations, tool
developers may choose to take this route for lower overhead at runtime
because the tool will not have to iterate over the entire set of variables
to find a specific one.
\end{rationale}

\begin{example}
\exindex{Tool information interface!listing names of all control variables}%
\Exindex{C}{Listing names of control variables}{MPI\_T\_cvar\_get\_info}
Querying and printing the names of all available control variables.

%%HEADER
%%LANG: C
%%ENDHEADER
\begin{lstlisting}[language={[MPI]C}]
#include <stdio.h>
#include <stdlib.h>
#include <mpi.h>

int main(int argc, char *argv[]) {
    int  i, err, num, namelen, bind, verbose, scope;
    int  threadsupport;
    char name[100];

    MPI_Datatype datatype;

    err=MPI_T_init_thread(MPI_THREAD_SINGLE, &threadsupport);
    if (err!=MPI_SUCCESS)
        return err;

    err=MPI_T_cvar_get_num(&num);
    if (err!=MPI_SUCCESS)
        return err;

    for (i=0; i<num; i++) {
        namelen=100;
        err=MPI_T_cvar_get_info(i, name, &namelen,
                                &verbose, &datatype, NULL,
                                NULL, NULL, /*no description */
                                &bind, &scope);
        if (err!=MPI_SUCCESS && err!=MPI_T_ERR_INVALID_INDEX)
            return err;
        printf("Var %i: %s\n", i, name);
    }

    err=MPI_T_finalize();
    if (err!=MPI_SUCCESS)
        return 1;
    else
        return 0;
}
\end{lstlisting}

\end{example}

\subsubsection{Handle Allocation and Deallocation}

Before reading or writing the value of a variable, a user must first
allocate a handle of type
\mpictypemain{MPI\_T\_cvar\_handle} for the
variable by binding it to an \MPI/ object (see also
Section~\ref{sec:mpit:assoc}).

\begin{rationale}
Handles used in the \MPI/ tool information interface are distinct from
handles used in the remaining parts of the \MPI/ standard because they must
be usable before \MPIisinitialized/ and after \MPIisfinalized/.
Further, accessing handles, in particular for performance variables, can be
time critical and having a separate handle space enables optimizations.
\end{rationale}

\cdeclindex{MPI\_T\_cvar\_handle}
\begin{funcdef}{MPI\_T\_CVAR\_HANDLE\_ALLOC}{(\mbox{cvar\_index}, \mbox{obj\_handle}, \mbox{handle}, \mbox{count})}
\funcarg{\IN}{cvar\_index}{index of control variable for which handle is to be allocated (index)}
\funcarg{\IN}{obj\_handle}{reference to a handle of the \mpi/ object to which this variable is supposed to be bound (pointer)}
\funcarg{\OUT}{handle}{allocated handle (handle)}
\funcarg{\OUT}{count}{number of elements used to represent this variable (integer)}
\end{funcdef}
\mpibindmain{MPI\_T\_cvar\_handle\_alloc}{(\gb{}int~cvar\_index, void~*obj\_handle, MPI\_T\_cvar\_handle~*handle, int~*count)}

This routine binds the control variable specified by the argument
\mpishortarg{index} to an \MPI/ object.  The object is passed in the
argument \mpiarg{obj\_handle} as an address to a local variable that
stores the object's handle.  The argument \mpiarg{obj\_handle} is ignored
if the \mpifunc{MPI\_T\_CVAR\_GET\_INFO} call for this control variable
returned \mpiconst{MPI\_T\_BIND\_NO\_OBJECT} in the argument
\mpishortarg{bind}.  The handle allocated to reference the variable is
returned in the argument \mpishortarg{handle}.  Upon successful return,
\mpishortarg{count} contains the number of elements (of the datatype
returned by a previous \mpifunc{MPI\_T\_CVAR\_GET\_INFO} call) used to
represent this variable.

\begin{users}
The \mpiarg{count} can be different based on the \MPI/ object to which the
control variable was bound.  For example, variables bound to communicators
could have a count that matches the size of the communicator.

It is not portable to pass references to predefined \MPI/ object handles,
such as \mpiconst{MPI\_COMM\_WORLD} to this routine, since their
implementation depends on the \MPI/ library.  Instead, such object handles
should be stored in a local variable and the address of this local variable
should be passed into \mpifunc{MPI\_T\_CVAR\_HANDLE\_ALLOC}.
\end{users}

The value of \mpiarg{cvar\_index} should be in the range from $0$ to
$\mpiarg{num\_cvar}-1$, where $\mpiarg{num\_cvar}$ is the number of available control
variables as determined from a prior call to
\mpifunc{MPI\_T\_CVAR\_GET\_NUM}.  The type of the \MPI/ object it
references must be consistent with the type returned in the \mpiarg{bind}
argument in a prior call to \mpifunc{MPI\_T\_CVAR\_GET\_INFO}.

\cdeclindex{MPI\_T\_cvar\_handle}
\begin{funcdef}{MPI\_T\_CVAR\_HANDLE\_FREE}{(\mbox{handle})}
\funcarg{\INOUT}{handle}{handle to be freed (handle)}
\end{funcdef}
\mpibindmain{MPI\_T\_cvar\_handle\_free}{(\gb{}MPI\_T\_cvar\_handle~*handle)}

When a handle is no longer needed, a user of the \MPI/ tool information
interface should call \mpifunc{MPI\_T\_CVAR\_HANDLE\_FREE} to free the
handle and the associated resources in the \MPI/ implementation.  On a
successful return, \MPI/ sets the handle to
\mpiconstmain{MPI\_T\_CVAR\_HANDLE\_NULL}.


\subsubsection{Control Variable Access Functions}

\cdeclindex{MPI\_T\_cvar\_handle}
\begin{funcdef}{MPI\_T\_CVAR\_READ}{(\mbox{handle}, \mbox{buf})}
\funcarg{\IN}{handle}{handle to the control variable to be read (handle)}
\funcarg{\OUT}{buf}{initial address of storage location for variable value (choice)}
\end{funcdef}
\mpibindmain{MPI\_T\_cvar\_read}{(\gb{}MPI\_T\_cvar\_handle~handle, void~*buf)}

This routine queries the value of a control variable identified by the
argument \mpiarg{handle} and stores the result in the buffer identified by
the parameter \mpiarg{buf}.  The user must ensure that the buffer is of the
appropriate size to hold the entire value of the control variable (based on
the returned datatype and count from prior corresponding calls to
\mpifunc{MPI\_T\_CVAR\_GET\_INFO} and
\mpifunc{MPI\_T\_CVAR\_HANDLE\_ALLOC}, respectively).

\cdeclindex{MPI\_T\_cvar\_handle}
\begin{funcdef}{MPI\_T\_CVAR\_WRITE}{(\mbox{handle}, \mbox{buf})}
\funcarg{\INOUT}{handle}{handle to the control variable to be written (handle)}
\funcarg{\IN}{buf}{initial address of storage location for variable value (choice)}
\end{funcdef}
\mpibindmain{MPI\_T\_cvar\_write}{(\gb{}MPI\_T\_cvar\_handle~handle, const~void~*buf)}

This routine sets the value of the control variable identified by the
argument \mpiarg{handle} to the data stored in the buffer identified by the
parameter \mpiarg{buf}.  The user must ensure that the buffer is of the
appropriate size to hold the entire value of the control variable (based on
the returned datatype and count from prior corresponding calls to
\mpifunc{MPI\_T\_CVAR\_GET\_INFO} and
\mpifunc{MPI\_T\_CVAR\_HANDLE\_ALLOC}, respectively).

If the variable has a global scope (as returned by a prior corresponding
\mpifunc{MPI\_T\_CVAR\_GET\_INFO} call), any write call to this variable
must be issued by the user in all connected (as defined in
Section~\ref{subsec:disconnect}) \MPI/ processes.  If the variable has
group scope, any write call to this variable must be issued by the user in
all \MPI/ processes in the group, which must be described by the \MPI/
implementation in the description returned by the call to \mpifunc{MPI\_T\_CVAR\_GET\_INFO}.

In both cases, the user must ensure that the writes in all participating \MPI/ processes are
consistent.  If the scope is either \mpiconst{MPI\_T\_SCOPE\_ALL\_EQ} or
\mpiconst{MPI\_T\_SCOPE\_GROUP\_EQ} this means that the variable in all
connected \MPI/ processes or \MPI/ processes of the group, respectively, must be set to the same value.

If it is not possible to change the variable at the time the call is made,
the function returns either \errormain{MPI\_T\_ERR\_CVAR\_SET\_NOT\_NOW}, if
there may be a later time at which the variable could be set, or
\errormain{MPI\_T\_ERR\_CVAR\_SET\_NEVER}, if the variable cannot be set for
the remainder of the application's execution.

\begin{example}
\exindex{Tool information interface!reading the value of a control variable}%
\Exindex{C}{Reading a control
  variable}{MPI\_T\_cvar\_handle\_alloc,MPI\_T\_cvar\_read,MPI\_T\_cvar\_handle\_free}%
Reading the value of a control variable.

%%HEADER
%%LANG: C
%%ENDHEADER
\begin{lstlisting}[language={[MPI]C}]
int getValue_int_comm(int index, MPI_Comm comm, int *val) {
    int err,count;
    MPI_T_cvar_handle handle;

    /* This example assumes that the variable index */
    /* can be bound to a communicator */

    err=MPI_T_cvar_handle_alloc(index, &comm, &handle, &count);
    if (err!=MPI_SUCCESS)
        return err;

    /* The following assumes that the variable is */
    /* represented by a single integer */

    err=MPI_T_cvar_read(handle,val);
    if (err!=MPI_SUCCESS)
        return err;

    err=MPI_T_cvar_handle_free(&handle);
    return err;
}
\end{lstlisting}
\end{example}


\subsection{Performance Variables}
\mpitermtitleindex{performance variable!tool information interface}
\mpitermindex{tool information interface!performance variable}
\label{sec:mpit:pvar}

The following section focuses on the ability to list and to query
performance variables provided by the \MPI/ implementation.  Performance
variables provide insight into \MPI/ implementation-specific internals and
can represent information such as the state of the \MPI/ implementation
(e.g., waiting blocked, receiving, not active), aggregated timing data for
submodules, or queue sizes and lengths.


\begin{rationale}
The interface for performance variables is separate from the interface for
control variables, since performance variables have different requirements
and parameters.  By keeping them separate, the interface provides cleaner
semantics and allows for more performance optimization opportunities.
\end{rationale}

Some performance variables and classes refer to \mpitermdefni{events}.
\mpitermdefindex{event!tool information interface}%
\mpitermdefindex{tool information interface!event}%
In general,
such events describe state transitions within software or hardware related
to the performance of an \MPI/ application. The events offered through the
callback-driven event-notification interface described in
Section~\ref{sec:mpit:events} also refer to such state transitions;
however, the set of state transitions referred to by performance variables and
events as described in Section~\ref{sec:mpit:events} may not be identical.

\subsubsection{Performance Variable Classes}
\label{sec:tools:mpit:classes}

Each performance variable is associated with a class that describes its
basic semantics, possible datatypes, basic behavior, its starting value,
whether it can overflow, and when and how an \MPI/ implementation can
change the variable's value.  The starting value is the value that is
assigned to the variable the first time that it is used or whenever it is
reset.

\begin{users}
If a performance variable belongs to a class that can overflow, it is up to
the user to protect against this overflow, e.g., by frequently reading and
resetting the variable value.
\end{users}

\begin{implementors}
\MPI/ implementations should use large enough datatypes for each
performance variable to avoid overflows under normal circumstances.
\end{implementors}

The classes are defined by the following constants:

\begin{sloppy}

\begin{description}

\item[\mpiconstmain{MPI\_T\_PVAR\_CLASS\_STATE}:]A performance variable in this
  class represents a set of discrete states.  Variables of this class are
  represented by \type{MPI\_INT} and can be set by the \MPI/
  implementation at any time.  Variables of this type should be described
  further using an enumeration, as discussed in Section~\ref{sec:types}.
  The starting value is the current state of the implementation at the time
  that the starting value is set.  \MPI/ implementations must ensure that
  variables of this class cannot overflow.

\item[\mpiconstmain{MPI\_T\_PVAR\_CLASS\_LEVEL}:]A performance variable in this
  class represents a value that describes the utilization level of a
  resource.  The value of a variable of this class can change at any time
  to match the current utilization level of the resource.  Values returned
  from variables in this class are nonnegative and represented by one of
  the following datatypes: \type{MPI\_UNSIGNED},
  \type{MPI\_UNSIGNED\_LONG}, \type{MPI\_UNSIGNED\_LONG\_LONG},
  \type{MPI\_DOUBLE}.  The starting value is the current utilization level
  of the resource at the time that the starting value is set.  \MPI/
  implementations must ensure that variables of this class cannot overflow.

\item[\mpiconstmain{MPI\_T\_PVAR\_CLASS\_SIZE}:]A performance variable in this
  class represents a value that is the size of a resource.  Values returned
  from variables in this class are nonnegative and represented by one of
  the following datatypes: \type{MPI\_UNSIGNED},
  \type{MPI\_UNSIGNED\_LONG}, \type{MPI\_UNSIGNED\_LONG\_LONG},
  \type{MPI\_DOUBLE}.  The starting value is the current size of the
  resource at the time that the starting value is set.  \MPI/
  implementations must ensure that variables of this class cannot overflow.

\item[\mpiconstmain{MPI\_T\_PVAR\_CLASS\_PERCENTAGE}:]The value of a performance
  variable in this class represents the percentage utilization of a finite
  resource.  The value of a variable of this class can change at any time
  to match the current utilization level of the resource.  It will be
  returned as an \type{MPI\_DOUBLE} datatype.  The value must always be
  between 0.0 (resource not used at all) and 1.0 (resource completely
  used).  The starting value is the current percentage utilization level of
  the resource at the time that the starting value is set.  \MPI/
  implementations must ensure that variables of this class cannot overflow.

\item[\mpiconstmain{MPI\_T\_PVAR\_CLASS\_HIGHWATERMARK}:]A performance variable
  in this class represents a value that describes the maximum observed
  utilization of a resource.  The value of a variable of this class is
  nonnegative and grows monotonically from the initialization or reset of
  the variable.  It can be represented by one of the following datatypes:
  \type{MPI\_UNSIGNED}, \type{MPI\_UNSIGNED\_LONG},
  \type{MPI\_UNSIGNED\_LONG\_LONG}, \type{MPI\_DOUBLE}.  The starting
  value is the current utilization level of the resource at the time that
  the variable is started or reset.  \MPI/ implementations must ensure that
  variables of this class cannot overflow.

\item[\mpiconstmain{MPI\_T\_PVAR\_CLASS\_LOWWATERMARK}:]A performance variable in
  this class represents a value that describes the minimum observed
  utilization of a resource.  The value of a variable of this class is
  nonnegative and decreases monotonically from the initialization or reset
  of the variable.  It can be represented by one of the following
  datatypes: \type{MPI\_UNSIGNED}, \type{MPI\_UNSIGNED\_LONG},
  \type{MPI\_UNSIGNED\_LONG\_LONG}, \type{MPI\_DOUBLE}.  The starting
  value is the current utilization level of the resource at the time that
  the variable is started or reset.  \MPI/ implementations must ensure that
  variables of this class cannot overflow.

\item[\mpiconstmain{MPI\_T\_PVAR\_CLASS\_COUNTER}:]A performance variable in this
  class counts the number of occurrences of a specific event (e.g., the
  number of memory allocations within an \MPI/ library).  The value of a
  variable of this class increases monotonically from the initialization or
  reset of the performance variable by one for each specific event that is
  observed.  Values must be nonnegative and represented by one of the
  following datatypes: \type{MPI\_UNSIGNED}, \type{MPI\_UNSIGNED\_LONG},
  \type{MPI\_UNSIGNED\_LONG\_LONG}.  The starting value for variables of
  this class is 0.  Variables of this class can overflow.

\item[\mpiconstmain{MPI\_T\_PVAR\_CLASS\_AGGREGATE}:]The value of a performance
  variable in this class is an aggregated value that represents a sum of
  arguments processed during a specific event (e.g., the amount of memory
  allocated by all memory allocations).  This class is similar to the
  counter class, but instead of counting individual events, the value can
  be incremented by arbitrary amounts.  The value of a variable of this
  class increases monotonically from the initialization or reset of the
  performance variable.  It must be nonnegative and represented by one of
  the following datatypes: \type{MPI\_UNSIGNED},
  \type{MPI\_UNSIGNED\_LONG}, \type{MPI\_UNSIGNED\_LONG\_LONG},
  \type{MPI\_DOUBLE}.  The starting value for variables of this class is
  0.  Variables of this class can overflow.

\item[\mpiconstmain{MPI\_T\_PVAR\_CLASS\_TIMER}:]The value of a performance
  variable in this class represents the aggregated time that the \MPI/
  implementation spends executing a particular event, type of event, or
  section of the \MPI/ library.  This class has the same basic semantics as
  \mpiconst{MPI\_T\_PVAR\_CLASS\_AGGREGATE}, but explicitly records a timing
  value.  The value of a variable of this class increases monotonically
  from the initialization or reset of the performance variable.  It must be
  nonnegative and represented by one of the following datatypes:
  \type{MPI\_UNSIGNED}, \type{MPI\_UNSIGNED\_LONG},
  \type{MPI\_UNSIGNED\_LONG\_LONG}, \type{MPI\_DOUBLE}.  The starting
  value for variables of this class is 0.  If the type \type{MPI\_DOUBLE}
  is used, the units that represent time in this datatype must match the
  units used by \mpifunc{MPI\_WTIME}.  Otherwise, the time units should be
  documented, e.g., in the description returned by
  \mpifunc{MPI\_T\_PVAR\_GET\_INFO}.  Variables of this class can overflow.

\item[\mpiconstmain{MPI\_T\_PVAR\_CLASS\_GENERIC}:]This class can be used to
  describe a variable that does not fit into any of the other classes.  For
  variables in this class, the starting value is variable-specific and
  implementation-defined.

\end{description}

\end{sloppy}

\subsubsection{Performance Variable Query Functions}

An \MPI/ implementation exports a set of $N$ performance variables through
the \MPI/ tool information interface.  If $N$ is zero, then the \MPI/
implementation does not export any performance variables; otherwise the
provided performance variables are indexed from $0$ to $N-1$.  This index
number is used in subsequent calls to identify the individual variables.

An \MPI/ implementation is allowed to increase the number of performance
variables during the execution of an \MPI/ application when new variables
become available through dynamic loading.  However, \MPI/ implementations
are not allowed to change the index of a performance variable or to delete
a variable once it has been added to the set.  When a variable becomes
inactive, e.g., through dynamic unloading, accessing its value should
return a corresponding return code.

The following function can be used to query the number of performance
variables, $\mpiarg{num\_pvar}$:

\begin{funcdef}{MPI\_T\_PVAR\_GET\_NUM}{(\mbox{num\_pvar})}
\funcarg{\OUT}{num\_pvar}{returns number of performance variables (integer)}
\end{funcdef}
\mpibindmain{MPI\_T\_pvar\_get\_num}{(\gb{}int~*num\_pvar)}

The function \mpifunc{MPI\_T\_PVAR\_GET\_INFO} provides access to
additional information for each variable.


\cdeclindex{MPI\_T\_enum}
\begin{funcdef}{MPI\_T\_PVAR\_GET\_INFO}{(\mbox{pvar\_index}, \mbox{name}, \mbox{name\_len}, \mbox{verbosity}, \mbox{var\_class}, \mbox{datatype}, \mbox{enumtype}, \mbox{desc}, \mbox{desc\_len}, \mbox{bind}, \mbox{readonly}, \mbox{continuous}, \mbox{atomic})}
\funcarg{\IN}{pvar\_index}{index of the performance variable to be queried between $0$ and $\mpiarg{num\_pvar}-1$ (integer)}
\funcarg{\OUT}{name}{buffer to return the string containing the name of the performance variable (string)}
\funcarg{\INOUT}{name\_len}{length of the string and/or buffer for \mpiarg{name} (integer)}
\funcarg{\OUT}{verbosity}{verbosity level of this variable (integer)}
\funcarg{\OUT}{var\_class}{class of performance variable (integer)}
\funcarg{\OUT}{datatype}{\mpi/ datatype of the information stored in the performance variable (handle)}
\funcarg{\OUT}{enumtype}{optional descriptor for enumeration information (handle)}
\funcarg{\OUT}{desc}{buffer to return the string containing a description of the performance variable (string)}
\funcarg{\INOUT}{desc\_len}{length of the string and/or buffer for \mpiarg{desc} (integer)}
\funcarg{\OUT}{bind}{type of \mpi/ object to which this variable must be bound (integer)}
\funcarg{\OUT}{readonly}{flag indicating whether the variable can be written/reset (integer)}
\funcarg{\OUT}{continuous}{flag indicating whether the variable can be started and stopped or is continuously active (integer)}
\funcarg{\OUT}{atomic}{flag indicating whether the variable can be atomically read and reset (integer)}
\end{funcdef}
\mpibindmain{MPI\_T\_pvar\_get\_info}{(\gb{}int~pvar\_index, char~*name, int~*name\_len, int~*verbosity, int~*var\_class, MPI\_Datatype~*datatype, MPI\_T\_enum~*enumtype, char~*desc, int~*desc\_len, int~*bind, int~*readonly, int~*continuous, int~*atomic)}

After a successful call to \mpifunc{MPI\_T\_PVAR\_GET\_INFO} for a
particular variable, subsequent calls to this routine that query
information about the same variable must return the same information.  An
\MPI/ implementation is not allowed to alter any of the returned values.

If any \OUT\ parameter to \mpifunc{MPI\_T\_PVAR\_GET\_INFO} is a
\code{NULL} pointer, the implementation will ignore the parameter and not
return a value for the parameter.

\textoutdesc{name}{the name of the performance variable} If completed
successfully, the routine is required to return a name of at least length
one.

The argument \mpiarg{verbosity} returns the verbosity level of the variable
(see Section~\ref{sec:mpit:verbose}).

The class of the performance variable is returned in the parameter
\mpiarg{var\_class}.  The class must be one of the constants defined in
Section~\ref{sec:tools:mpit:classes}.

The combination of the name and the class of the performance variable must
be unique with respect to all other names for  performance variables used
by the \MPI/ implementation.

\begin{implementors}
Groups of variables that belong closely together, but have different
classes, can have the same name.  This choice is useful, e.g., to refer to
multiple variables that describe a single resource (like the level, the
total size, as well as high- and low-water marks).
\end{implementors}

The argument \mpiarg{datatype} returns the \MPI/ datatype that is used to
represent the performance variable.

If the variable is of type \type{MPI\_INT}, \MPI/ can optionally specify
an enumeration for the values represented by this variable and return it in
\mpiarg{enumtype}.  In this case, \MPI/ returns an enumeration identifier,
which can then be used to gather more information as described in
Section~\ref{sec:types}.  Otherwise, \mpiarg{enumtype} is set to
\mpiconst{MPI\_T\_ENUM\_NULL}.  If the datatype is not \type{MPI\_INT} or the
argument \mpiarg{enumtype} is the null pointer, no enumeration type is
returned.

Returning a description is optional.  If an \mpi/ implementation does not
return a description, the first character for \mpiarg{desc} must be set to
the null character and \mpiarg{desc\_len} must be set to one at the return
from this function.

The parameter \mpiarg{bind} returns the type of the \MPI/ object to which
the variable must be bound or the value \mpiconst{MPI\_T\_BIND\_NO\_OBJECT}
(see Section~\ref{sec:mpit:assoc}).

Upon return, the argument \mpiarg{readonly} is set to zero if the variable
can be written or reset by the user.  It is set to one if the variable can
only be read.

Upon return, the argument \mpiarg{continuous} is set to zero if the
variable can be started and stopped by the user, i.e., it is possible for
the user to control if and when the value of a variable is updated.  It is
set to one if the variable is always active and cannot be controlled by the
user.

Upon return, the argument \mpiarg{atomic} is set to zero if the variable
cannot be read and reset atomically.  Only variables for which the call
sets \mpiarg{atomic} to one can be used in a call to
\mpifunc{MPI\_T\_PVAR\_READRESET}.

If a performance variable has an equivalent name and has the same class
across connected \MPI/ processes, the following \OUT\  parameters must be
identical: \mpiarg{verbosity}, \mpiarg{varclass}, \mpiarg{datatype},
\mpiarg{enumtype}, \mpiarg{bind}, \mpiarg{readonly}, \mpiarg{continuous},
and \mpiarg{atomic}.  The returned description must be equivalent.

\begin{funcdef}{MPI\_T\_PVAR\_GET\_INDEX}{(\mbox{name}, \mbox{var\_class}, \mbox{pvar\_index})}
\funcarg{\IN}{name}{the name of the performance variable (string)}
\funcarg{\IN}{var\_class}{the class of the performance variable (integer)}
\funcarg{\OUT}{pvar\_index}{the index of the performance variable (integer)}
\end{funcdef}
\mpibindmain{MPI\_T\_pvar\_get\_index}{(\gb{}const~char~*name, int~var\_class, int~*pvar\_index)}

\mpifunc{MPI\_T\_PVAR\_GET\_INDEX} is a function for retrieving the index
of a performance variable given a known variable name and class.  The
\mpiarg{name} and \mpiarg{var\_class} parameters are provided by the
caller, and \mpiarg{pvar\_index} is returned by the \MPI/ implementation.
The \mpiarg{name} parameter is a string terminated with a null character.

This routine returns \error{MPI\_SUCCESS} on success and returns
\error{MPI\_T\_ERR\_INVALID\_NAME} if \mpiarg{name} does not match the name
of any performance variable of the specified \mpiarg{var\_class} provided
by the implementation at the time of the call.

\begin{rationale}
This routine is provided to enable fast retrieval of performance variables
by a tool, assuming it knows the name of the variable for which it is
looking.  The number of variables exposed by the implementation can change
over time, so it is not possible for the tool to simply iterate over the
list of variables once at initialization.  Although using \MPI/
implementation specific variable names is not portable across \MPI/
implementations, tool developers may choose to take this route for lower
overhead at runtime because the tool will not have to iterate over the
entire set of variables to find a specific one.
\end{rationale}

\subsubsection{Performance Experiment Sessions}
\label{sec:mpit:sessions}

Within a single program, multiple components can use the \MPI/ tool
information interface.  To avoid collisions with respect to accesses to
performance variables, users of the \MPI/ tool information interface must
first create a performance experiment session.  Subsequent calls that
access performance variables can then be made within the context of this
performance experiment session.   Starting, stopping, reading, writing, or
resetting a variable in one performance experiment session shall not
influence whether a variable is started, stopped, read, written, or reset
in another performance experiment session.

\cdeclindex{MPI\_T\_pvar\_session}
\begin{funcdef}{MPI\_T\_PVAR\_SESSION\_CREATE}{(\mbox{pe\_session})}
\funcarg{\OUT}{pe\_session}{identifier of performance experiment session (handle)}
\end{funcdef}
\mpibindmain{MPI\_T\_pvar\_session\_create}{(\gb{}MPI\_T\_pvar\_session~*pe\_session)}

This call creates a new performance experiment session for accessing performance variables and
returns a handle for this performance experiment session in the argument \mpiarg{pe\_session} of type
\mpictypemain{MPI\_T\_pvar\_session}.

\cdeclindex{MPI\_T\_pvar\_session}
\begin{funcdef}{MPI\_T\_PVAR\_SESSION\_FREE}{(\mbox{pe\_session})}
\funcarg{\INOUT}{pe\_session}{identifier of performance experiment session (handle)}
\end{funcdef}
\mpibindmain{MPI\_T\_pvar\_session\_free}{(\gb{}MPI\_T\_pvar\_session~*pe\_session)}

This call frees an existing performance experiment session.  Calls to the \MPI/ tool information
interface  can no longer be made within the context of a performance experiment session after it
is freed.  On a successful return, \MPI/ sets the performance experiment session identifier to
\mpiconstmain{MPI\_T\_PVAR\_SESSION\_NULL}.


\subsubsection{Handle Allocation and Deallocation}

Before using a performance variable, a user must first allocate a handle of
type \mpictypemain{MPI\_T\_pvar\_handle} for
the variable by binding it to an \MPI/ object (see also
Section~\ref{sec:mpit:assoc}).

\cdeclindex{MPI\_T\_pvar\_handle}
\cdeclindex{MPI\_T\_pvar\_session}
\begin{funcdef}{MPI\_T\_PVAR\_HANDLE\_ALLOC}{(\mbox{pe\_session}, \mbox{pvar\_index}, \mbox{obj\_handle}, \mbox{handle}, \mbox{count})}
\funcarg{\INOUT}{pe\_session}{identifier of performance experiment session (handle)}
\funcarg{\IN}{pvar\_index}{index of performance variable for which handle is to be allocated (integer)}
\funcarg{\IN}{obj\_handle}{reference to a handle of the \mpi/ object to which this variable is supposed to be bound (pointer)}
\funcarg{\OUT}{handle}{allocated handle (handle)}
\funcarg{\OUT}{count}{number of elements used to represent this variable (integer)}
\end{funcdef}
\mpibindmain{MPI\_T\_pvar\_handle\_alloc}{(\gb{}MPI\_T\_pvar\_session~pe\_session, int~pvar\_index, void~*obj\_handle, MPI\_T\_pvar\_handle~*handle, int~*count)}

This routine binds the performance variable specified by the argument
\mpiarg{index} to an \MPI/ object  in the performance experiment session identified by the
parameter \mpishortarg{pe\_session}.  The object is passed in the argument
\mpiarg{obj\_handle} as an address to a local variable that stores the
object's handle.  The argument \mpiarg{obj\_handle} is ignored if the
\mpifunc{MPI\_T\_PVAR\_GET\_INFO} call for this performance variable
returned \mpiconst{MPI\_T\_BIND\_NO\_OBJECT} in the argument \mpiarg{bind}.
The handle allocated to reference the variable is returned in the argument
\mpiarg{handle}.  Upon successful return, \mpiarg{count} contains the
number of elements (of the datatype returned by a previous
\mpifunc{MPI\_T\_PVAR\_GET\_INFO} call) used to represent this variable.

\begin{users}
The \mpiarg{count} can be different based on the \MPI/ object to which the
performance variable was bound.  For example, variables bound to
communicators could have a count that matches the size of the communicator.

It is not portable to pass references to predefined \MPI/ object handles,
such as \mpiconst{MPI\_COMM\_WORLD}, to this routine, since their
implementation depends on the \MPI/ library.  Instead, such an object
handle should be stored in a local variable and the address of this local
variable should be passed into \mpifunc{MPI\_T\_PVAR\_HANDLE\_ALLOC}.
\end{users}

The value of index should be in the range from $0$ to $\mpiarg{num\_pvar}-1$,
where $\mpishortarg{num\_pvar}$ is the number of available performance
variables as determined from a prior call to
\mpifunc{MPI\_T\_PVAR\_GET\_NUM}.  The type of the \MPI/ object it
references must be consistent with the type returned in the \mpiarg{bind}
argument in a prior call to \mpifunc{MPI\_T\_PVAR\_GET\_INFO}.

For all routines in the rest of this section that take both \mpiarg{handle}
and \mpiarg{pe\_session} as \IN\ or \INOUT\ arguments, if the \mpiarg{handle}
argument passed in is not associated with the \mpiarg{pe\_session} argument,
\errormain{MPI\_T\_ERR\_INVALID\_HANDLE} is returned.

\cdeclindex{MPI\_T\_pvar\_handle}%
\cdeclindex{MPI\_T\_pvar\_session}%
\begin{funcdef}{MPI\_T\_PVAR\_HANDLE\_FREE}{(\mbox{pe\_session}, \mbox{handle})}
\funcarg{\INOUT}{pe\_session}{identifier of performance experiment session (handle)}
\funcarg{\INOUT}{handle}{handle to be freed (handle)}
\end{funcdef}
\mpibindmain{MPI\_T\_pvar\_handle\_free}{(\gb{}MPI\_T\_pvar\_session~pe\_session, MPI\_T\_pvar\_handle~*handle)}

When a handle is no longer needed, a user of the \MPI/ tool information
interface should call \mpifunc{MPI\_T\_PVAR\_HANDLE\_FREE} to free the
handle in the performance experiment session identified by the parameter \mpiarg{pe\_session} and the
associated resources in the \MPI/  implementation.  On a successful return,
\MPI/ sets the handle to \mpiconstmain{MPI\_T\_PVAR\_HANDLE\_NULL}.


\subsubsection{Starting and Stopping of Performance Variables}

Performance variables that have the continuous flag set during the query
procedure are continuously updated once a handle has been allocated.
Such variables may be queried at any time, but they cannot be started or
stopped by the user.  All other variables are in a stopped state after
their handle has been allocated; their values are not updated until they
have been started by the user.

\cdeclindex{MPI\_T\_pvar\_handle}
\cdeclindex{MPI\_T\_pvar\_session} 
\begin{funcdef}{MPI\_T\_PVAR\_START}{(\mbox{pe\_session}, \mbox{handle})}
\funcarg{\IN}{pe\_session}{identifier of performance experiment session (handle)}
\funcarg{\INOUT}{handle}{handle of a performance variable (handle)}
\end{funcdef}
\mpibindmain{MPI\_T\_pvar\_start}{(\gb{}MPI\_T\_pvar\_session~pe\_session, MPI\_T\_pvar\_handle~handle)}

This functions starts the performance variable with the handle identified
by the parameter \mpiarg{handle} in the performance experiment session identified by the parameter
\mpiarg{pe\_session}.

If the constant \mpiconstmain{MPI\_T\_PVAR\_ALL\_HANDLES} is passed in
\mpiarg{handle}, the \MPI/ implementation attempts to start all variables
within the performance experiment session identified by the parameter \mpishortarg{pe\_session} for
which handles have been allocated.  In this case, the routine returns
\error{MPI\_SUCCESS} if all variables are started successfully (even if
there are no noncontinuous variables to be started), otherwise
\errormain{MPI\_T\_ERR\_PVAR\_NO\_STARTSTOP} is returned.  Continuous variables
and variables that are already started are ignored when
\mpiconst{MPI\_T\_PVAR\_ALL\_HANDLES} is specified.

\cdeclindex{MPI\_T\_pvar\_handle} 
\cdeclindex{MPI\_T\_pvar\_session}
\begin{funcdef}{MPI\_T\_PVAR\_STOP}{(\mbox{pe\_session}, \mbox{handle})}
\funcarg{\IN}{pe\_session}{identifier of performance experiment session (handle)}
\funcarg{\INOUT}{handle}{handle of a performance variable (handle)}
\end{funcdef}
\mpibindmain{MPI\_T\_pvar\_stop}{(\gb{}MPI\_T\_pvar\_session~pe\_session, MPI\_T\_pvar\_handle~handle)}

This functions stops the performance variable with the handle identified by
the parameter \mpiarg{handle} in the performance experiment session identified by the parameter
\mpiarg{pe\_session}.

If the constant \mpiconst{MPI\_T\_PVAR\_ALL\_HANDLES} is passed in
\mpiarg{handle}, the \MPI/ implementation attempts to stop all variables
within the performance experiment session identified by the parameter \mpishortarg{pe\_session} for
which handles have been allocated.  In this case, the routine returns
\error{MPI\_SUCCESS} if all variables are stopped successfully (even if
there are no noncontinuous variables to be stopped), otherwise
\error{MPI\_T\_ERR\_PVAR\_NO\_STARTSTOP} is returned.  Continuous variables
and variables that are already stopped are ignored when
\mpiconst{MPI\_T\_PVAR\_ALL\_HANDLES} is specified.


\subsubsection{Performance Variable Access Functions}

\cdeclindex{MPI\_T\_pvar\_handle} 
\cdeclindex{MPI\_T\_pvar\_session}
\begin{funcdef}{MPI\_T\_PVAR\_READ}{(\mbox{pe\_session}, \mbox{handle}, \mbox{buf})}
\funcarg{\IN}{pe\_session}{identifier of performance experiment session (handle)}
\funcarg{\IN}{handle}{handle of a performance variable (handle)}
\funcarg{\OUT}{buf}{initial address of storage location for variable value (choice)}
\end{funcdef}
\mpibindmain{MPI\_T\_pvar\_read}{(\gb{}MPI\_T\_pvar\_session~pe\_session, MPI\_T\_pvar\_handle~handle, void~*buf)}

The \mpifunc{MPI\_T\_PVAR\_READ} call queries the value of the performance
variable with the handle \mpiarg{handle} in the performance experiment session identified by the
parameter \mpiarg{pe\_session} and stores the result in the buffer identified
by the parameter \mpiarg{buf}.  The user must ensure that the
buffer is of the appropriate size to hold the entire value of the
performance variable (based on the datatype and count returned by the
corresponding previous calls to \mpifunc{MPI\_T\_PVAR\_GET\_INFO} and
\mpifunc{MPI\_T\_PVAR\_HANDLE\_ALLOC}, respectively).

The constant \mpiconst{MPI\_T\_PVAR\_ALL\_HANDLES} cannot be used as an
argument for the function \mpifunc{MPI\_T\_PVAR\_READ}.

\cdeclindex{MPI\_T\_pvar\_handle} 
\cdeclindex{MPI\_T\_pvar\_session}
\begin{funcdef}{MPI\_T\_PVAR\_WRITE}{(\mbox{pe\_session}, \mbox{handle}, \mbox{buf})}
\funcarg{\IN}{pe\_session}{identifier of performance experiment session (handle)}
\funcarg{\INOUT}{handle}{handle of a performance variable (handle)}
\funcarg{\IN}{buf}{initial address of storage location for variable value (choice)}
\end{funcdef}
\mpibindmain{MPI\_T\_pvar\_write}{(\gb{}MPI\_T\_pvar\_session~pe\_session, MPI\_T\_pvar\_handle~handle, const~void~*buf)}

The \mpifunc{MPI\_T\_PVAR\_WRITE} call attempts to write the value of the
performance variable with the handle identified by the parameter
\mpiarg{handle} in the performance experiment session identified by the parameter
\mpiarg{pe\_session}.  The value to be written is passed in the buffer
identified by the parameter \mpiarg{buf}.  The user must ensure that the
buffer is of the appropriate size to hold the entire value of the
performance variable (based on the datatype and count returned by the
corresponding previous calls to \mpifunc{MPI\_T\_PVAR\_GET\_INFO} and
\mpifunc{MPI\_T\_PVAR\_HANDLE\_ALLOC}, respectively).

If it is not possible to change the variable, the function returns
\errormain{MPI\_T\_ERR\_PVAR\_NO\_WRITE}.

The constant \mpiconst{MPI\_T\_PVAR\_ALL\_HANDLES} cannot be used as an
argument for the  function \mpifunc{MPI\_T\_PVAR\_WRITE}.

\cdeclindex{MPI\_T\_pvar\_handle} 
\cdeclindex{MPI\_T\_pvar\_session}
\begin{funcdef}{MPI\_T\_PVAR\_RESET}{(\mbox{pe\_session}, \mbox{handle})}
\funcarg{\IN}{pe\_session}{identifier of performance experiment session (handle)}
\funcarg{\INOUT}{handle}{handle of a performance variable (handle)}
\end{funcdef}
\mpibindmain{MPI\_T\_pvar\_reset}{(\gb{}MPI\_T\_pvar\_session~pe\_session, MPI\_T\_pvar\_handle~handle)}

The \mpifunc{MPI\_T\_PVAR\_RESET} call sets the performance variable with
the handle identified by the parameter \mpiarg{handle} to its starting
value specified in Section~\ref{sec:tools:mpit:classes}.  If it is not
possible to change the variable, the function returns
\error{MPI\_T\_ERR\_PVAR\_NO\_WRITE}.

If the constant \mpiconst{MPI\_T\_PVAR\_ALL\_HANDLES} is passed in
\mpiarg{handle}, the \MPI/ implementation attempts to reset all variables
within the performance experiment session identified by the parameter \mpiarg{pe\_session} for which
handles have been allocated.  In this case, the routine returns
\error{MPI\_SUCCESS} if all variables are reset successfully (even if there
are no valid handles or all are read-only), otherwise
\error{MPI\_T\_ERR\_PVAR\_NO\_WRITE} is returned.  Read-only variables are
ignored when \mpiconst{MPI\_T\_PVAR\_ALL\_HANDLES} is specified.


\cdeclindex{MPI\_T\_pvar\_handle} 
\cdeclindex{MPI\_T\_pvar\_session}
\begin{funcdef}{MPI\_T\_PVAR\_READRESET}{(\mbox{pe\_session}, \mbox{handle}, \mbox{buf})}
\funcarg{\IN}{pe\_session}{identifier of performance experiment session (handle)}
\funcarg{\INOUT}{handle}{handle of a performance variable (handle)}
\funcarg{\OUT}{buf}{initial address of storage location for variable value (choice)}
\end{funcdef}
\mpibindmain{MPI\_T\_pvar\_readreset}{(\gb{}MPI\_T\_pvar\_session~pe\_session, MPI\_T\_pvar\_handle~handle, void~*buf)}

This call atomically combines the functionality of
\mpifunc{MPI\_T\_PVAR\_READ} and \mpifunc{MPI\_T\_PVAR\_RESET} with the
same semantics as if these two calls were called separately.  If the
variable cannot be read and reset atomically, this routine returns
\errormain{MPI\_T\_ERR\_PVAR\_NO\_ATOMIC}.

The constant \mpiconst{MPI\_T\_PVAR\_ALL\_HANDLES} cannot be used as an
argument for the  function \mpifunc{MPI\_T\_PVAR\_READRESET}.

\begin{implementors}
Sampling-based tools rely on the ability to call the \MPI/ tool information
interface, in particular routines to start, stop, read, write, and reset
performance variables, from any program context, including asynchronous
contexts such as signal handlers.  \MPI/ implementations should strive, if
possible in their particular environment, to enable these usage scenarios
for all or a subset of the routines mentioned above.  If implementing only
a subset, the read, write, and reset routines are typically the most
critical for sampling-based tools.  An \MPI/ implementation should clearly
document any restrictions on the program contexts in which the \MPI/ tool
information interface can be used.  Restrictions might include guaranteeing
usage outside of all signals or outside a specific set of signals.  Any
restrictions could be documented, for example, through the description
returned by \mpifunc{MPI\_T\_PVAR\_GET\_INFO}.
\end{implementors}

\begin{rationale}
All routines to read, to write or to reset performance variables require
the performance experiment session argument.  This requirement keeps the interface consistent and
allows the use of \mpiconst{MPI\_T\_PVAR\_ALL\_HANDLES}  where appropriate.
Further, this opens up additional performance optimizations for the
implementation of handles.
\end{rationale}


\begin{example}
\exindex{Tool information interface!basic usage of performance variables}%
Detecting receives with long unexpected message queues.
\Exindex{C}{Basic usage of performance variables}{MPI\_T\_init\_thread,MPI\_T\_pvar\_get\_info,MPI\_T\_pvar\_session\_create,MPI\_T\_pvar\_handle\_alloc,MPI\_T\_pvar\_start,MPI\_T\_pvar\_read,MPI\_T\_pvar\_handle\_free,MPI\_T\_pvar\_handle\_free,MPI\_T\_finalize}

The following example shows a sample tool to identify receive operations
that occur during times with long message queues.  This examples assumes
that the \MPI/ implementation exports a variable with the name
``\code{MPI\_T\_UMQ\_LENGTH}'' to represent the current length of the
unexpected message queue.  The tool is implemented as a \mpiskipfunc{PMPI}
tool using the
\MPI/ profiling interface.

The tool consists of three parts: (1) the initialization (by intercepting
the call to \mpifunc{MPI\_INIT}), (2) the test for long unexpected message
queues (by intercepting calls to \mpifunc{MPI\_RECV}), and (3) the clean-up
phase (by intercepting the call to \mpifunc{MPI\_FINALIZE}).  To capture
all receives, the example would have to be extended to have similar
wrappers for all receive operations.

\paraheading{Part 1---Initialization:}
During initialization, the tool searches for the variable and, once the
right index is found, allocates a performance experiment session and a handle for the variable
with the found index, and starts the performance variable.

%%HEADER
%%LANG: C
%%ENDHEADER
\begin{lstlisting}[language={[MPI]C},basicstyle=\lstsmall]
#include <stdio.h>
#include <stdlib.h>
#include <string.h>
#include <assert.h>
#include <mpi.h>

/* Global variables for the tool */
static MPI_T_pvar_session pe_session;
static MPI_T_pvar_handle handle;

int MPI_Init(int *argc, char ***argv ) {
    int  err, num, i, index, namelen, verbosity;
    int  var_class, bind, threadsup;
    int  readonly, continuous, atomic, count;
    char name[18];

    MPI_Comm     comm;
    MPI_Datatype datatype;
    MPI_T_enum   enumtype;

    err=PMPI_Init(argc, argv);
    if (err!=MPI_SUCCESS)
        return err;

    err=PMPI_T_init_thread(MPI_THREAD_SINGLE, &threadsup);
    if (err!=MPI_SUCCESS)
        return err;

    err=PMPI_T_pvar_get_num(&num);
    if (err!=MPI_SUCCESS)
        return err;

    index=-1;
    i=0;
    while ((i<num) && (index<0) && (err==MPI_SUCCESS)) {
        /* Pass a buffer that is at least one character longer than */
        /* the name of the variable being searched for to avoid */
        /* finding variables that have a name that has a prefix */
        /* equal to the name of the variable being searched. */
        namelen=18;
        err=PMPI_T_pvar_get_info(i, name, &namelen, &verbosity,
                                 &var_class, &datatype, &enumtype,
                                 NULL, NULL, &bind,&readonly,
                                 &continuous, &atomic);
        if (strcmp(name,"MPI_T_UMQ_LENGTH")==0) index=i;
        i++;
    }
    if (err!=MPI_SUCCESS)
        return err;

    /* this could be handled in a more flexible way for a generic tool */
    assert(index>=0);
    assert(var_class==MPI_T_PVAR_CLASS_LEVEL);
    assert(datatype==MPI_INT);
    assert(bind==MPI_T_BIND_MPI_COMM);

    /* Create a session */
    err=PMPI_T_pvar_session_create(&pe_session);
    if (err!=MPI_SUCCESS) return err;

    /* Get a handle and bind to MPI_COMM_WORLD */
    comm=MPI_COMM_WORLD;
    err=PMPI_T_pvar_handle_alloc(pe_session, index, &comm, &handle,
                                 &count);
    if (err!=MPI_SUCCESS) return err;

    /* this could be handled in a more flexible way for a generic tool */
    assert(count==1);

    /* Start variable */
    err=PMPI_T_pvar_start(pe_session, handle);
    if (err!=MPI_SUCCESS) return err;

    return MPI_SUCCESS;
}
\end{lstlisting}

\paraheading{Part 2---Testing the Queue Lengths During Receives:}
During every receive operation, the tool reads the unexpected queue length
through the matching performance variable and compares it against a
predefined threshold.

%%HEADER
%%LANG: C
%%DECL: MPI_T_pvar_session pe_session; MPI_T_pvar_handle handle;
%%ENDHEADER
\begin{lstlisting}[language={[MPI]C},basicstyle=\lstsmall]
#define THRESHOLD 5

int MPI_Recv(void *buf, int count, MPI_Datatype datatype, int source,
             int tag, MPI_Comm comm, MPI_Status *status)
{
        int value, err;

        if (comm==MPI_COMM_WORLD) {
                err=PMPI_T_pvar_read(pe_session, handle, &value);
                if ((err==MPI_SUCCESS) && (value>THRESHOLD))
                {
                        /* tool identified receive called with long UMQ */
                        /* execute tool functionality, */
                        /* e.g., gather and print call stack */
                }
        }

        return PMPI_Recv(buf, count, datatype, source, tag, comm, status);
}
\end{lstlisting}

\paraheading{Part 3---Termination:}
In the wrapper for \mpifunc{MPI\_FINALIZE}, the \MPI/ tool information
interface is finalized.

%%HEADER
%%LANG: C
%%DECL: MPI_T_pvar_session pe_session; MPI_T_pvar_handle handle;
%%SUBST: return:if (err) return err;return
%%ENDHEADER
\begin{lstlisting}[language={[MPI]C}]
int MPI_Finalize(void)
{
    int err;

    err=PMPI_T_pvar_handle_free(pe_session, &handle);
    err=PMPI_T_pvar_session_free(&pe_session);
    err=PMPI_T_finalize();
    return PMPI_Finalize();
}
\end{lstlisting}


\end{example}

\subsection{Events}
\mpitermtitleindex{event!tool information interface}
\mpitermindex{tool information interface!event}
\label{sec:mpit:events}

During the execution of an \MPI/ application, the \MPI/ implementation can
raise \mpitermni{events}\mpitermdefindex{event!tool information interface}
\mpitermdefindex{tool information interface!event} of a specific type to inform the user of a state change
in the implementation.  \mpitermdefni{Event types}\mpitermdefindex{event type!tool information interface}\mpitermdefindex{tool information interface!event type}
describe specific state changes
within the \MPI/ implementation.  In comparison to aggregate performance
variables, events provide per-instance information on such state changes.
The \MPI/ implementation is said to \mpitermdef{raise an event} when it invokes a
callback function previously registered by the user for the corresponding event
type.  Each callback invocation for a specific event instance has a
timestamp associated with it, which can be queried by the user, describing
the time when the event was observed by the implementation. This decouples
the observation of the state change from the communication of this
information to the user.  A timestamp in this context is a count of clock
ticks elapsed since some time in the past and represented as a variable of
type \mpictype{MPI\_Count}.

\subsubsection{Event Sources}
\label{sec:tools:sources}

As a means to manage multiple state changes to be observed concurrently by
different parts of the software and hardware system, the event interface of
the \MPI/ Tool Information Interface uses the concept of \emph{sources}.  A
source in this context is a concept describing the logical entity raising
the event. A source may or may not directly represent a concrete part of
the software or hardware system. This concept is used primarily to describe
partial ordering of events across different components where total ordering
cannot necessarily be determined or is too costly to enforce. 

The following function can be used to query the number of event sources, \mpiarg{num\_sources}:

\begin{funcdef}{MPI\_T\_SOURCE\_GET\_NUM}{(\mbox{num\_sources})}
\funcarg{\OUT}{num\_sources}{returns number of event sources (integer)}
\end{funcdef}
\mpibindmain{MPI\_T\_source\_get\_num}{(\gb{}int~*num\_sources)}

The number of available event sources can be queried with a call to
\mpifunc{MPI\_T\_SOURCE\_GET\_NUM}.  An \MPI/ implementation is allowed to
increase the number of sources during the execution of an \MPI/ process.
However, \MPI/ implementations are not allowed to change the index of an
event source or to delete an event source once it has been made visible to
the user (e.g., if new event sources become available via dynamic loading
of additional components in the \MPI/ implementation).

\cdeclindex{MPI\_T\_source\_order}
\begin{funcdef}{MPI\_T\_SOURCE\_GET\_INFO}{(\mbox{source\_index}, \mbox{name}, \mbox{name\_len}, \mbox{desc}, \mbox{desc\_len}, \mbox{ordering}, \mbox{ticks\_per\_second}, \mbox{max\_ticks}, \mbox{info})}
\funcarg{\IN}{source\_index}{index of the source to be queried between $0$ and $\mpiarg{num\_sources}-1$ (integer)}
\funcarg{\OUT}{name}{buffer to return the string containing the name of the source (string)}
\funcarg{\INOUT}{name\_len}{length of the string and/or buffer for \mpiarg{name} (integer)}
\funcarg{\OUT}{desc}{buffer to return the string containing the description of the source (string)}
\funcarg{\INOUT}{desc\_len}{length of the string and/or buffer for \mpiarg{desc} (integer)}
\funcarg{\OUT}{ordering}{flag indicating chronological ordering guarantees given by the source (integer)}
\funcarg{\OUT}{ticks\_per\_second}{the number of ticks per second for the timer of this source (integer)}
\funcarg{\OUT}{max\_ticks}{the maximum count of ticks reported by this source before overflow occurs (integer)}
\funcarg{\OUT}{info}{optional info object (handle)}
\end{funcdef}
\mpibindmain{MPI\_T\_source\_get\_info}{(\gb{}int~source\_index, char~*name, int~*name\_len, char~*desc, int~*desc\_len, MPI\_T\_source\_order~*ordering, MPI\_Count~*ticks\_per\_second, MPI\_Count~*max\_ticks, MPI\_Info~*info)}

A call to \mpifunc{MPI\_T\_SOURCE\_GET\_INFO} returns additional
information on the source identified by the \mpiarg{source\_index}
argument.

\textoutdesc{name}{the name of the source}

\textoutdesc{desc}{the description of the source}

The \mpiarg{ordering} argument is of type \mpictypemain{MPI\_T\_source\_order} and returns whether event callbacks of this
source will be invoked in chronological order, i.e., the timestamps reported by
\mpifunc{MPI\_T\_EVENT\_GET\_TIMESTAMP} of subsequent events of the same source are
monotonically increasing. The value of \mpiarg{ordering} can be
\mpiconstmain{MPI\_T\_SOURCE\_ORDERED} or \mpiconstmain{MPI\_T\_SOURCE\_UNORDERED}.

The \mpiarg{ticks\_per\_seconds} argument returns the number of ticks elapsed in
one second for the timer used for the specific source.

The \mpiarg{max\_ticks} argument returns the largest number of ticks
reported by this source as a timestamp before the value overflows.

\begin{users}
    As the size of \mpictype{MPI\_Count} is defined in relation to the types
    \mpictype{MPI\_Aint} and \mpictype{MPI\_Offset}, the effective size of
    \mpictype{MPI\_Count} may lead to overflows of the timestamp values
    reported.  Users can use the argument \mpiarg{max\_ticks} to
    mitigate resulting problems.
\end{users}


\textoutoptinfo{info}{this source}

\begin{funcdef}{MPI\_T\_SOURCE\_GET\_TIMESTAMP}{(\mbox{source\_index}, \mbox{timestamp})}
\funcarg{\IN}{source\_index}{index of the source (integer)}
\funcarg{\OUT}{timestamp}{current timestamp from specified source (integer)}
\end{funcdef}
\mpibindmain{MPI\_T\_source\_get\_timestamp}{(\gb{}int~source\_index, MPI\_Count~*timestamp)}

To enable proper query of a reference timestamp for a specific source, a
user can obtain a current timestamp using
\mpifunc{MPI\_T\_SOURCE\_GET\_TIMESTAMP}.  The argument
\mpiarg{source\_index} identifies the index of the source to query.  The
call returns \error{MPI\_SUCCESS} and a current timestamp in the argument
\mpiarg{timestamp} if the source supports ad-hoc generation of timestamps.
The call returns \error{MPI\_T\_ERR\_INVALID\_INDEX} if the index does not
identify a valid source.  The call returns
\errormain{MPI\_T\_ERR\_NOT\_SUPPORTED} if the source does not support the
ad-hoc generation of timestamps.


\subsubsection{Callback Safety Requirements}
\cdeclmainindex{MPI\_T\_cb\_safety}

The actions a user is allowed to perform inside a callback function may vary
with its execution context.  As the user has no control over the execution
context of specific callback function invocations, \MPI/ provides a way to
communicate this information using callback safety levels.

\begin{table}[t]
\caption{Hierarchy of safety requirement levels for event callback routines}
\label{table:tools:mpit:cbsafety}
\begin{center}
\begin{tabular}{|l|}
\hline
\textbf{Safety Requirement} \\
\hline
\mpiconst{MPI\_T\_CB\_REQUIRE\_NONE} \\
\mpiconst{MPI\_T\_CB\_REQUIRE\_MPI\_RESTRICTED} \\
\mpiconst{MPI\_T\_CB\_REQUIRE\_THREAD\_SAFE} \\
\mpiconst{MPI\_T\_CB\_REQUIRE\_ASYNC\_SIGNAL\_SAFE} \\
\hline
\end{tabular}
\end{center}
\end{table}

Table~\ref{table:tools:mpit:cbsafety} provides the hierarchy of
callback safety requirements levels within user-defined callback functions.
The \MPI/ implementation provides the
safety requirement as an argument to the callback when it is invoked.

The level of \mpiconstmain{MPI\_T\_CB\_REQUIRE\_NONE} is the lowest
level and does not impose any restrictions on the callback function.  

The level of \mpiconstmain{MPI\_T\_CB\_REQUIRE\_MPI\_RESTRICTED} restricts the
set of \MPI/ functions that can be called from inside the callback to 
all functions with the prefix \mpiskipfunc{MPI\_T} as well as \mpifunc{MPI\_WTICK} 
and \mpifunc{MPI\_WTIME}.

\begin{users}
    While some \MPI/ functions are safe to be called inside a callback
    function used in the \MPI/ tool information interface---which
    may in some implementations be issued from asynchronous
    contexts such as signal handlers---this does not imply that those \MPI/
    functions are generally safe to be called in asynchronous contexts such as
    signal handlers.
\end{users}

The level of \mpiconstmain{MPI\_T\_CB\_REQUIRE\_THREAD\_SAFE} includes all the
limitations of \mpiconst{MPI\_T\_CB\_REQUIRE\_MPI\_RESTRICTED} and
additionally requires the callback to be reentrant and thread-safe\mpitermindex{threads!thread-safe}.  This
means the callback must allow its execution to be
interrupted by or happen concurrently with any other callback including
itself.

The level of \mpiconstmain{MPI\_T\_CB\_REQUIRE\_ASYNC\_SIGNAL\_SAFE} includes all
the limitations of \mpiconst{MPI\_T\_CB\_REQUIRE\_THREAD\_SAFE} and
additionally requires the callback to meet the safety requirements 
needed to support invocations from asynchronous contexts, such
as signal handlers.

\begin{users}
    It is always safe to assume the highest restrictions for a callback
    invocation (i.e., \mpiconst{MPI\_T\_CB\_REQUIRE\_ASYNC\_SIGNAL\_SAFE}).
    By evaluating the specific requirements at runtime, a tool may obtain
    more freedom of action within the callback.
\end{users}

\begin{implementors}
    A high-quality implementation will strive to set callback safety
    requirements to the most permissive level for a given callback
    invocation.
\end{implementors}

All functions with the prefix \mpiskipfunc{MPI\_T}, except those listed in
Table~\ref{table:mpit:mpi_restricted}, may return the return code
\errormain{MPI\_T\_ERR\_NOT\_ACCESSIBLE}
to indicate that the user may not access this function at this time.
The functions (and their respective \mpiskipfunc{PMPI} versions) listed in
Table~\ref{table:mpit:mpi_restricted} are exceptions to this rule and shall
not return \error{MPI\_T\_ERR\_NOT\_ACCESSIBLE}.

\begin{table}[t]
    \caption{\label{table:mpit:mpi_restricted}List of \MPI/ functions that when called from within a callback function may not return \error{MPI\_T\_ERR\_NOT\_ACCESSIBLE}}
\begin{center}
    \begin{tabular}{l}
        \mpifunc{MPI\_T\_EVENT\_COPY} \\
        \mpifunc{MPI\_T\_EVENT\_GET\_SOURCE} \\
        \mpifunc{MPI\_T\_EVENT\_GET\_TIMESTAMP} \\
        \mpifunc{MPI\_T\_EVENT\_READ} \\
        \mpifunc{MPI\_T\_PVAR\_READ} \\
        \mpifunc{MPI\_T\_PVAR\_READRESET} \\
        \mpifunc{MPI\_T\_PVAR\_RESET} \\
        \mpifunc{MPI\_T\_PVAR\_START} \\
        \mpifunc{MPI\_T\_PVAR\_STOP} \\
        \mpifunc{MPI\_T\_PVAR\_WRITE} \\
        \mpifunc{MPI\_T\_SOURCE\_GET\_TIMESTAMP} \\
    \end{tabular}
\end{center}
\end{table}

\begin{rationale}
    A call may be implemented in a way that is not safe for all execution
    contexts of a callback function, e.g., inside a signal handler. An \MPI/
    implementation therefore needs a way to communicate its inability to perform
    a certain action due to the execution context of a callback invocation.
\end{rationale}

\begin{implementors}
    A high-quality implementation shall not return \error{MPI\_T\_ERR\_NOT\_ACCESSIBLE}
    except where absolutely necessary.
\end{implementors}

\begin{users}
    Users intercepting calls into the \MPI/ tool information interface using
    the \mpiskipfunc{PMPI} interface must ensure that the safety requirements for the calling
    context are met.  This means that users may have to implement the
    wrapper with the highest safety level used by the \MPI/ implementation.
\end{users}


\subsubsection{Event Type Query Functions}

An \MPI/ implementation exports a set of $N$ event types through the \MPI/
tool information interface. If $N$ is zero, then the \MPI/ implementation
does not export any event types; otherwise, the provided event types
are indexed from $0$ to $N-1$. This index number is used in subsequent
calls to identify a specific event type.

An \MPI/ implementation is allowed to increase the number of event types
during the execution of an \MPI/ process.
However, \MPI/ implementations are not
allowed to change the index of an event type or to delete an event type
once it has been made visible to the user (e.g., if new event types
become available via dynamic loading of additional components in the
\MPI/ implementation).

The following function can be used to query the number of event types, \mpiarg{num\_events}:

\begin{funcdef}{MPI\_T\_EVENT\_GET\_NUM}{(\mbox{num\_events})}
\funcarg{\OUT}{num\_events}{returns number of event types (integer)}
\end{funcdef}
\mpibindmain{MPI\_T\_event\_get\_num}{(\gb{}int~*num\_events)}

The function \mpifunc{MPI\_T\_EVENT\_GET\_INFO} provides access to
additional information about a specific event type.

\cdeclindex{MPI\_T\_enum}
\begin{funcdef}{MPI\_T\_EVENT\_GET\_INFO}{(\mbox{event\_index}, \mbox{name}, \mbox{name\_len}, \mbox{verbosity}, \mbox{array\_of\_datatypes}, \mbox{array\_of\_displacements}, \mbox{num\_elements}, \mbox{enumtype}, \mbox{info}, \mbox{desc}, \mbox{desc\_len}, \mbox{bind})}
\funcarg{\IN}{event\_index}{index of the event type to be queried between $0$ and $\mpiarg{num\_events}-1$ (integer)}
\funcarg{\OUT}{name}{buffer to return the string containing the name of the event type (string)}
\funcarg{\INOUT}{name\_len}{length of the string and/or buffer for \mpiarg{name} (integer)}
\funcarg{\OUT}{verbosity}{verbosity level of this event type (integer)}
\funcarg{\OUT}{array\_of\_datatypes}{array of \mpi/ basic datatypes used to encode the event data (array of handles)}
\funcarg{\OUT}{array\_of\_displacements}{array of byte displacements of the elements in the event buffer (array of nonnegative integers)}
\funcarg{\INOUT}{num\_elements}{length of \mpiarg{array\_of\_datatypes} and \mpiarg{array\_of\_displacements} arrays (nonnegative integer)}
\funcarg{\OUT}{enumtype}{optional descriptor for enumeration information (handle)}
\funcarg{\OUT}{info}{optional info object (handle)}
\funcarg{\OUT}{desc}{buffer to return the string containing a description of the event type (string)}
\funcarg{\INOUT}{desc\_len}{length of the string and/or buffer for \mpiarg{desc} (integer)}
\funcarg{\OUT}{bind}{type of \mpi/ object to which an event of this type must be bound (integer)}
\end{funcdef}
\mpibindmain{MPI\_T\_event\_get\_info}{(\gb{}int~event\_index, char~*name, int~*name\_len, int~*verbosity, MPI\_Datatype~array\_of\_datatypes[], MPI\_Aint~array\_of\_displacements[], int~*num\_elements, MPI\_T\_enum~*enumtype, MPI\_Info~*info, char~*desc, int~*desc\_len, int~*bind)}

After a successful call to \mpifunc{MPI\_T\_EVENT\_GET\_INFO} for a
particular event type, subsequent calls to this routine that query
information about the same event type must return the same information.  If
any \INOUT\ or \OUT\ argument to
\mpifunc{MPI\_T\_EVENT\_GET\_INFO} is a \code{NULL} pointer, the
implementation will ignore the argument and not return a value for the
specific argument.

\textoutdesc{name}{the name of the event type} If completed successfully,
the routine is required to return a name of at least length one.  The name
of the event type must be unique with respect to all other names for event
types used by the \MPI/ implementation.

The argument \mpiarg{verbosity} returns the verbosity level of the event
type (see Section~\ref{sec:mpit:verbose}).

The argument \mpiarg{array\_of\_datatypes} returns an array of \MPI/
datatype handles that describe the elements returned for an instance
of the event type with index \mpiarg{event\_index}.
The event data can either be queried element by element with 
\mpifunc{MPI\_T\_EVENT\_READ} or copied
into a contiguous event buffer with \mpifunc{MPI\_T\_EVENT\_COPY}. 
For the latter case, the 
argument \mpiarg{array\_of\_displacements} returns an 
array of byte displacements in the event buffer in ascending order starting with zero.

The user is responsible
for the memory allocation for the \mpiarg{array\_of\_datatypes} and
\mpiarg{array\_of\_displacements} arrays. The number of elements
in each array is supplied by the user in \mpiarg{num\_elements}.
If the number of elements used by the event type is larger
than the value of \mpiarg{num\_elements} provided by the user, the number
of datatype handles and displacements returned in the corresponding arrays
is truncated to the value of \mpiarg{num\_elements} passed in by the user.
If the user passes the \code{NULL} pointer for \mpiarg{array\_of\_datatypes} 
or \mpiarg{array\_of\_displacements}, the respective arguments are ignored.
Unless the user passes the \code{NULL} pointer for \mpiarg{num\_elements}, 
the function returns the number of elements required for this event type. 
If the user passes the \code{NULL} pointer for \mpiarg{num\_elements}, 
the arguments \mpiarg{num\_elements},
\mpiarg{array\_of\_datatypes}, and \mpiarg{array\_of\_displacements} are
ignored.


\MPI/ can optionally return an enumeration identifier in the
\mpiarg{enumtype} argument, describing the individual elements in the
\mpiarg{array\_of\_datatypes} argument. Otherwise, \mpiarg{enumtype} is set
to \mpiconst{MPI\_T\_ENUM\_NULL}. If the argument to \mpiarg{enumtype}
provided by the user is the \code{NULL} pointer, no enumeration
type is returned.

\textoutoptinfo{info}{a registration handle for this event type}

\textoutdesc{desc}{the description of the event type} Returning a
description is optional.  If an \mpi/ implementation does not return a
description, the first character for \mpiarg{desc} must be set to the null
character and \mpiarg{desc\_len} must be set to one at the return from this
function.

The parameter \mpiarg{bind} returns the type of the \MPI/ object to which
the event type must be bound or the value \mpiconst{MPI\_T\_BIND\_NO\_OBJECT}
(see Section~\ref{sec:mpit:assoc}).

If an event type has an equivalent name across connected \MPI/ processes, the
following \OUT\ parameters must be identical: \mpiarg{verbosity},
\mpiarg{array\_of\_datatypes}, \mpiarg{num\_elements}, \mpiarg{enumtype},
and \mpiarg{bind}.  The returned description must be equivalent.
As the argument \mpiarg{array\_of\_displacements} is process dependent, it
may differ across connected \MPI/ processes.

This routine returns \error{MPI\_SUCCESS} on success and returns
\error{MPI\_T\_ERR\_INVALID\_INDEX} if \mpiarg{event\_index} does not match
a valid event type index provided by the implementation at the time of the
call.

\begin{funcdef}{MPI\_T\_EVENT\_GET\_INDEX}{(\mbox{name}, \mbox{event\_index})}
\funcarg{\IN}{name}{name of the event type (string)}
\funcarg{\OUT}{event\_index}{index of the event type (integer)}
\end{funcdef}
\mpibindmain{MPI\_T\_event\_get\_index}{(\gb{}const~char~*name, int~*event\_index)}

\mpifunc{MPI\_T\_EVENT\_GET\_INDEX} returns the index of an event type
identified by a known event type name.  The \mpiarg{name} parameter is
provided by the caller, and \mpiarg{event\_index} is returned by the \MPI/
implementation. The \mpiarg{name} parameter is a string terminated with a
null character.

This routine returns \error{MPI\_SUCCESS} on success and returns
\error{MPI\_T\_ERR\_INVALID\_NAME} if \mpiarg{name} does not match the name
of any event type provided by the implementation at the time of the call.

\begin{rationale}
This routine is provided to enable fast retrieval of an event index by a 
tool, assuming it knows the name of the event type for which it is looking.
The number of event types exposed by the implementation can change over
time, so it is not possible for the tool to simply iterate over the list of
event types once at initialization.  Although using \MPI/ implementation
specific event type names is not portable across \MPI/ implementations, tool
developers may choose to take this route for lower overhead at runtime
because the tool will not have to iterate over the entire set of event types 
to find a specific one.
\end{rationale}


\subsubsection{Handle Allocation and Deallocation}
\label{sec:tools:mpit:eventalloc}

Before the \MPI/ implementation calls a callback function on the occurrence of
a specific event, the user needs to register a callback function to be called for
that event type and obtain a handle of type
\mpictypemain{MPI\_T\_event\_registration}.

\cdeclindex{MPI\_T\_event\_registration}
\begin{funcdef}{MPI\_T\_EVENT\_HANDLE\_ALLOC}{(\mbox{event\_index}, \mbox{obj\_handle}, \mbox{info}, \mbox{event\_registration})}
\funcarg{\IN}{event\_index}{index of event type for which the registration handle is to be allocated (integer)}
\funcarg{\IN}{obj\_handle}{reference to a handle of the \mpi/ object to which this event is supposed to be bound (pointer)}
\funcarg{\IN}{info}{info object (handle)}
\funcarg{\OUT}{event\_registration}{event registration (handle)}
\end{funcdef}
\mpibindmain{MPI\_T\_event\_handle\_alloc}{(\gb{}int~event\_index, void~*obj\_handle, MPI\_Info~info, MPI\_T\_event\_registration~*event\_registration)}

\mpifunc{MPI\_T\_EVENT\_HANDLE\_ALLOC} creates a \mpitermdef{registration handle}
for the event type identified by \mpiarg{event\_index}. Furthermore,
if required by the event type, the registration handle is bound to the
object referred to by the argument \mpiarg{obj\_handle}.  The argument
\mpiarg{obj\_handle} is ignored if the \mpifunc{MPI\_T\_EVENT\_GET\_INFO}
call for this event type returned \mpiconst{MPI\_T\_BIND\_NO\_OBJECT} in the
argument \mpishortarg{bind}.  The user can pass hints for the handle
allocation to the \MPI/ implementation via the \mpiarg{info} argument.  The
allocated event-registration handle is returned in the argument
\mpiarg{event\_registration}.

\cdeclindex{MPI\_T\_event\_registration}
\begin{funcdef}{MPI\_T\_EVENT\_HANDLE\_SET\_INFO}{(\mbox{event\_registration}, \mbox{info})}
\funcarg{\INOUT}{event\_registration}{event registration (handle)}
\funcarg{\IN}{info}{info object (handle)}
\end{funcdef}
\mpibindmain{MPI\_T\_event\_handle\_set\_info}{(\gb{}MPI\_T\_event\_registration~event\_registration, MPI\_Info~info)}

\mpifunc{MPI\_T\_EVENT\_HANDLE\_SET\_INFO} updates the hints of the
event-registration handle associated with \mpiarg{event\_registration} using the
hints provided in \mpiarg{info}. A call to this procedure has no effect on previously
set or defaulted hints that are not specified by \mpiarg{info}. It also has
no effect on previously set or defaulted hints that are specified by info,
but are ignored by the \MPI/ implementation in this call to
\mpifunc{MPI\_T\_EVENT\_HANDLE\_SET\_INFO}.


\begin{users}
    Some info items that an implementation can use when it creates an
    event-registration handle cannot easily be changed once the
    registration handle is created. Thus, an implementation may ignore
    hints issued in this call that it would have accepted in a handle
    allocation call. An implementation may also be unable to update certain
    info hints in a call to \mpifunc{MPI\_T\_EVENT\_HANDLE\_SET\_INFO}.
    \mpifunc{MPI\_T\_EVENT\_HANDLE\_GET\_INFO} can be used to determine
    whether info changes were ignored by the implementation.
\end{users}


\cdeclindex{MPI\_T\_event\_registration}
\begin{funcdef}{MPI\_T\_EVENT\_HANDLE\_GET\_INFO}{(\mbox{event\_registration}, \mbox{info\_used})}
\funcarg{\IN}{event\_registration}{event registration (handle)}
\funcarg{\OUT}{info\_used}{info object (handle)}
\end{funcdef}
\mpibindmain{MPI\_T\_event\_handle\_get\_info}{(\gb{}MPI\_T\_event\_registration~event\_registration, MPI\_Info~*info\_used)}

\mpifunc{MPI\_T\_EVENT\_HANDLE\_GET\_INFO} returns a new info object
containing the hints of the event-registration handle associated with
\mpiarg{event\_registration}.  The current setting of all hints related to
this registration handle is returned in \mpiarg{info\_used}.  An \MPI/
implementation is required to return all hints that are supported by the
implementation and have default values specified; any user-supplied hints
that were not ignored by the implementation; and any additional hints that
were set by the implementation.  If no such hints exist, a handle to a
newly created info object is returned that contains no key/value pairs.
The user is responsible for freeing \mpiarg{info\_used} via
\mpifunc{MPI\_INFO\_FREE}.


\begin{funcdef}{MPI\_T\_EVENT\_REGISTER\_CALLBACK}{(\mbox{event\_registration}, \mbox{cb\_safety}, \mbox{info}, \mbox{user\_data}, \mbox{event\_cb\_function})}
\funcarg{\INOUT}{event\_registration}{event registration (handle)}
\funcarg{\IN}{cb\_safety}{maximum callback safety level (integer)}
\funcarg{\IN}{info}{info object (handle)}
\funcarg{\IN}{user\_data}{pointer to a user-controlled buffer}
\funcarg{\IN}{event\_cb\_function}{pointer to user-defined callback function (function)}
\end{funcdef}
\mpibindmain{MPI\_T\_event\_register\_callback}{(\gb{}MPI\_T\_event\_registration~event\_registration, MPI\_T\_cb\_safety~cb\_safety, MPI\_Info~info, void~*user\_data, MPI\_T\_event\_cb\_function~event\_cb\_function)}
\cdeclindex{MPI\_T\_cb\_safety}
\cdeclindex{MPI\_T\_event\_registration}
\mpicallbackindex{MPI\_T\_event\_cb\_function}

\mpifunc{MPI\_T\_EVENT\_REGISTER\_CALLBACK} associates a user-defined
function pointed to by \mpiarg{event\_cb\_function} with an
allocated event-registration handle.
The maximum callback safety level supported
by the callback function is passed in the argument \mpiarg{cb\_safety}.  
The safety levels are defined in Table~\ref{table:tools:mpit:cbsafety}.
A user can
register multiple callback functions for a given event-registration handle, potentially
specifying one for each callback safety level. 
Registering a callback function for a specific callback safety level
overwrites any previously-registered callback function pointer and info
object associated with the event
registration for the specific callback safety level. If
\mpiarg{event\_cb\_function} is the \code{NULL} pointer, an existing association
of a callback function for that callback safety level is removed.

When an event is triggered, the implementation will select from all registered callbacks
the callback with the lowest safety level valid in the context in which the callback
is invoked.  In situations where the required callback
safety level exceeds the highest level for which a callback function is 
registered for a given registration handle, the event instance is dropped.  

At callback invocation time, the implementation passes the pointer to a user-defined memory region specified during
callback registration with the argument \mpiarg{user\_data}.

The user can pass hints for
the registration of the specified callback function to the \MPI/ implementation via the \mpiarg{info}
argument.

\begin{users}
    As event instances can be raised as soon as the registration handle is
    associated with the first callback function, the callback function with
    the highest callback safety guarantees should be registered before any
    further registrations for lower callback safety guarantees, to avoid
    dropped events due to insufficient callback safety guarantees.
\end{users}

The callback function passed to \mpifunc{MPI\_T\_EVENT\_REGISTER\_CALLBACK}
in the argument \mpiarg{event\_cb\_function} needs to have the following
type:

\mpitypedefbindvoidmain{MPI\_T\_event\_cb\_function}{(\gb{}MPI\_T\_event\_instance~event\_instance, MPI\_T\_event\_registration~event\_registration, MPI\_T\_cb\_safety~cb\_safety, void~*user\_data)}

The argument \mpiarg{event\_instance} corresponds to a handle for the
opaque event-instance object of type 
\mpictypemain{MPI\_T\_event\_instance}. 
This handle is only valid inside the
corresponding invocation of the function to which it is passed.  The
argument \mpiarg{event\_registration} corresponds to the event-registration handle
returned by \mpifunc{MPI\_T\_EVENT\_HANDLE\_ALLOC} for the user function to
the same event type and bound object combination.  The handle can be used
to identify the specific event registration information, such as event type
and bound object, or even to deallocate the handle from within the callback
invocation.
The argument \mpiarg{cb\_safety} describes the safety requirements the
callback function must fulfill in the current invocation.  The argument
\mpiarg{user\_data} is the pointer to user-allocated memory that was
passed to the \MPI/ implementation during callback registration.

\cdeclindex{MPI\_T\_cb\_safety}
\cdeclindex{MPI\_T\_event\_registration}
\begin{funcdef}{MPI\_T\_EVENT\_CALLBACK\_SET\_INFO}{(\mbox{event\_registration}, \mbox{cb\_safety}, \mbox{info})}
\funcarg{\INOUT}{event\_registration}{event registration (handle)}
\funcarg{\IN}{cb\_safety}{callback safety level (integer)}
\funcarg{\IN}{info}{info object (handle)}
\end{funcdef}
\mpibindmain{MPI\_T\_event\_callback\_set\_info}{(\gb{}MPI\_T\_event\_registration~event\_registration, MPI\_T\_cb\_safety~cb\_safety, MPI\_Info~info)}

\mpifunc{MPI\_T\_EVENT\_CALLBACK\_SET\_INFO} updates the hints of the
callback function registered for the callback safety level specified by
\mpiarg{cb\_safety} of the event-registration handle associated with
\mpiarg{event\_registration} using the hints provided in \mpiarg{info}.
A call to this procedure has no effect on previously set or defaulted hints that are
not specified by \mpiarg{info}. It also has no effect on previously set or
defaulted hints that are specified by info, but are ignored by the \MPI/
implementation in this call to
\mpifunc{MPI\_T\_EVENT\_CALLBACK\_SET\_INFO}.

\cdeclindex{MPI\_T\_cb\_safety}
\cdeclindex{MPI\_T\_event\_registration}
\begin{funcdef}{MPI\_T\_EVENT\_CALLBACK\_GET\_INFO}{(\mbox{event\_registration}, \mbox{cb\_safety}, \mbox{info\_used})}
\funcarg{\IN}{event\_registration}{event registration (handle)}
\funcarg{\IN}{cb\_safety}{callback safety level (integer)}
\funcarg{\OUT}{info\_used}{info object (handle)}
\end{funcdef}
\mpibindmain{MPI\_T\_event\_callback\_get\_info}{(\gb{}MPI\_T\_event\_registration~event\_registration, MPI\_T\_cb\_safety~cb\_safety, MPI\_Info~*info\_used)}

\mpifunc{MPI\_T\_EVENT\_CALLBACK\_GET\_INFO} returns a new info object
containing the hints of the callback function registered for the callback
safety level specified by \mpiarg{cb\_safety} of the event-registration
handle associated with \mpiarg{event\_registration}.  The current set
of all hints related to this callback safety level of the
event-registration handle is returned in \mpiarg{info\_used}.
An \MPI/ implementation is required to return all hints that are supported
by the implementation and have default values specified, any user-supplied
hints that were not ignored by the implementation, and any additional hints
that were set by the implementation.  If no such hints exist, a handle to a
newly created info object is returned that contains no key/value pairs.
The user is responsible for freeing \mpiarg{info\_used} via
\mpifunc{MPI\_INFO\_FREE}.


To stop the \MPI/ implementation from raising events for a specific
registration, a user needs to free the corresponding event-registration
handle.

\begin{funcdef}{MPI\_T\_EVENT\_HANDLE\_FREE}{(\mbox{event\_registration}, \mbox{user\_data}, \mbox{free\_cb\_function})}
\funcarg{\INOUT}{event\_registration}{event registration (handle)}
\funcarg{\IN}{user\_data}{pointer to a user-controlled buffer}
\funcarg{\IN}{free\_cb\_function}{pointer to user-defined callback function (function)}
\end{funcdef}
\mpibindmain{MPI\_T\_event\_handle\_free}{(\gb{}MPI\_T\_event\_registration~event\_registration, void~*user\_data, MPI\_T\_event\_free\_cb\_function~free\_cb\_function)}

\mpifunc{MPI\_T\_EVENT\_HANDLE\_FREE} returns \error{MPI\_SUCCESS} when
deallocation of the handle was initiated successfully and returns
\error{MPI\_T\_ERR\_INVALID\_HANDLE} if \mpiarg{event\_registration} does
not match a valid allocated event-registration handle at the time of the call.
The callback function \mpiarg{free\_cb\_function} is called by the \MPI/
implementation, when it is able to guarantee that no further event
instances for the corresponding event-registration handle will be raised.  If
the pointer to \mpiarg{free\_cb\_function} is the \code{NULL} pointer, no
user function is invoked after successful deallocation of the event
registration handle. The pointer to user-controlled memory provided in the
\mpiarg{user\_data} argument will be passed to the function provided in the
\mpiarg{free\_cb\_function} on invocation.

\begin{users}
    A free-callback function associated with a registration handle should
    always be prepared to postpone any pending actions, should the provided
    callback safety requirements exceed those required by the pending
    actions.
\end{users}


The callback function passed to \mpifunc{MPI\_T\_EVENT\_HANDLE\_FREE} in
the argument \mpiarg{free\_cb\_function} needs to have the following
type:

\mpitypedefbindvoidmain{MPI\_T\_event\_free\_cb\_function}{(\gb{}MPI\_T\_event\_registration~event\_registration, MPI\_T\_cb\_safety~cb\_safety, void~*user\_data)}

\subsubsection{Handling Dropped Events}
\label{sec:tools:mpit:droppedevents}

Events may occur at times when the \MPI/ implementation cannot invoke
the user function corresponding to a matching event handle.
An implementation is allowed to buffer such events and delay
the callback invocation. If an event occurs at times when
the corresponding callback function cannot be called and the
corresponding data cannot be buffered,
or no callback function meeting the required callback safety level is registered,
the event data may be
dropped. 
To discover such data loss, the user can set a handler
function for a specific event-registration handle.

\begin{funcdef}{MPI\_T\_EVENT\_SET\_DROPPED\_HANDLER}{(\mbox{event\_registration}, \mbox{dropped\_cb\_function})}
\funcarg{\INOUT}{event\_registration}{valid event registration (handle)}
\funcarg{\IN}{dropped\_cb\_function}{pointer to user-defined callback function (function)}
\end{funcdef}
\mpibindmain{MPI\_T\_event\_set\_dropped\_handler}{(\gb{}MPI\_T\_event\_registration~event\_registration, MPI\_T\_event\_dropped\_cb\_function~dropped\_cb\_function)}
\cdeclindex{MPI\_T\_event\_registration}
\mpicallbackindex{MPI\_T\_event\_dropped\_cb\_function}

\mpifunc{MPI\_T\_EVENT\_SET\_DROPPED\_HANDLER} registers the function 
\mpiarg{dropped\_cb\_function} to be called by the \MPI/
implementation when event information is dropped for the registration handle
specified in \mpiarg{event\_registration}. 
Subsequent calls to \mpifunc{MPI\_T\_EVENT\_SET\_DROPPED\_HANDLER} 
with the same registration handle will replace previously-registered callback functions 
for that registration handle.
If the pointer to
\mpiarg{dropped\_cb\_function} is the \code{NULL} pointer,
no data loss is recorded or reported until a new valid callback function is registered.


\begin{users}
    The invocation of the dropped handler callback function may not necessarily occur 
    close to the time the event was actually lost.
\end{users}


The callback function passed to \mpifunc{MPI\_T\_EVENT\_SET\_DROPPED\_HANDLER} in
the argument \mpiarg{dropped\_cb\_function} needs to have the following
type:

\mpitypedefbindvoidmain{MPI\_T\_event\_dropped\_cb\_function}{(\gb{}MPI\_Count~count, MPI\_T\_event\_registration~event\_registration, int~source\_index, MPI\_T\_cb\_safety~cb\_safety, void~*user\_data)}

The argument \mpiarg{event\_registration} corresponds to the event
registration handle to which the dropped data corresponds. The argument
\mpiarg{count} provides a best effort estimation of the number of
invocations to a registered event callback corresponding to
\mpiarg{event\_registration} that were not executed since the registration
of the dropped-callback handler or the last invocation of a registered 
dropped-callback handler.  If the number of dropped events observed by the
implementation exceeds the limit of \mpiargshort{count}, an
implementation shall set \mpiargshort{count} to the maximum possible value for the type of \mpiargshort{count}.  The
\mpiarg{source\_index} provides the index of the source that dropped the
corresponding event information. The argument \mpiargshort{cb\_safety} describes
the safety requirements the callback function must fulfill in the current
invocation.  The possible values for \mpiargshort{cb\_safety} are described in
Table~\ref{table:tools:mpit:cbsafety}. The argument \mpiarg{user\_data} is
the pointer to user-allocated memory that was passed to the \MPI/
implementation during callback registration.  If no event callback is
registered for safety requirement levels that an implementation uses to invoke
the dropped handler callback function for a specific event, the corresponding
dropped handler callback function will not be invoked.

\begin{users}
    A callback function for dropped events associated with a registration
    handle should always be prepared to postpone any pending actions,
    should the provided callback safety requirements exceed those required
    by the pending actions.
\end{users}

\begin{implementors}
    A high-quality implementation should strive to find a good balance between
    timely notification, completeness of information, and the freedom of action
    for a tool when invoking the callback function for dropped events associated
    with a registration handle.
\end{implementors}

If dropped event notifications have been observed for a specific source since the last
event notification of that source, the corresponding dropped handler
callback function must be called before other events are raised for
that source. This means in a sequence of five events E1 to E5 from the same
source, where E3 and E4 were dropped, any handler function set through
\mpifunc{MPI\_T\_EVENT\_SET\_DROPPED\_HANDLER} for event-registration
handles associated with E3 or E4 must be called before E5 is raised.


\subsubsection{Reading Event Data}

In event callbacks, the parameter \mpiarg{event\_instance} provides
access to the per-instance event data, i.e., the data encoded by the
specific event type for this instance.  The user can obtain event data as
well as event meta data, such as a time stamp and the source, by providing
this handle to the respective query functions. The event-instance handle is
invalid beyond the scope of the current invocation of the callback function
to which it is provided.

The callback function argument \mpiarg{event\_registration} 
identifies the registration handle that was used to register the callback function.

The callback function  argument \mpiarg{cb\_safety} indicates the requirements for the
specific callback invocation. The value is one of the safety requirements levels
described in Table~\ref{table:tools:mpit:cbsafety}.
The argument \mpiarg{user\_data} passes the pointer provided by the
user during callback registration back to the function call.

\begin{users}
Depending on the registered event and usage of \MPI/ by the application, a
callback function may be invoked with high frequency. Users should
therefore strive to minimize the amount of work done inside callback
functions. Furthermore, the time spent in a callback function may influence
the capability of an implementation to buffer events; long execution
times may lead to an increased number of dropped events.
\end{users}

\MPI/ provides the following function calls to access data of a specific event
instance and its corresponding meta data (such as its time and source).

\cdeclindex{MPI\_T\_event\_instance}
\begin{funcdef}{MPI\_T\_EVENT\_READ}{(\mbox{event\_instance}, \mbox{element\_index}, \mbox{buffer})}
\funcarg{\IN}{event\_instance}{event-instance handle provided to the callback function (handle)}
\funcarg{\IN}{element\_index}{index into the array of datatypes of the item to be queried (integer)}
\funcarg{\OUT}{buffer}{pointer to a memory location to store the item data (choice)}
\end{funcdef}
\mpibindmain{MPI\_T\_event\_read}{(\gb{}MPI\_T\_event\_instance~event\_instance, int~element\_index, void~*buffer)}

\mpifunc{MPI\_T\_EVENT\_READ} allows users to copy one element of the event
data to a user-specified buffer at a time.

The \mpiarg{event\_instance} argument identifies the event instance to
query. It is erroneous to provide any other event-instance handle to the
call than the one passed by the \MPI/ implementation to the callback
function in which the data is read.  The \mpishortarg{buffer} argument must point
to a memory location the \MPI/ implementation can copy the element
of the event data to identified by \mpiarg{element\_index}.

\cdeclindex{MPI\_T\_event\_instance}
\begin{funcdef}{MPI\_T\_EVENT\_COPY}{(\mbox{event\_instance}, \mbox{buffer})}
\funcarg{\IN}{event\_instance}{event instance provided to the callback function (handle)}
\funcarg{\OUT}{buffer}{user-allocated buffer for event data (choice)}
\end{funcdef}
\mpibindmain{MPI\_T\_event\_copy}{(\gb{}MPI\_T\_event\_instance~event\_instance, void~*buffer)}

\mpifunc{MPI\_T\_EVENT\_COPY} copies the event data as a whole into the
user-provided \mpiarg{buffer}. The user must ensure that the buffer is of
at least the size of the extent of the event type, which can be computed from the type and displacement information returned by the
corresponding call to \mpifunc{MPI\_T\_EVENT\_GET\_INFO}. The data may
include padding bytes between individual elements of the event data in the
buffer. A user can reconstruct the location and size of the data contained
in the buffer through the information returned by
\mpifunc{MPI\_T\_EVENT\_GET\_INFO}.

\begin{implementors}
    An implementation should strive to use an appropriately compact
    representation when copying event instance data to a user buffer via
    \mpifunc{MPI\_T\_EVENT\_COPY} to reduce the amount of memory required for
    the user buffer.
\end{implementors}

\subsubsection{Reading Event Meta Data}
\label{sec:tools:mpit:events:metadata}

Additional to the specific event data encoded by each event type,
supplemental information available across all event types can be queried.

\cdeclindex{MPI\_T\_event\_instance}
\begin{funcdef}{MPI\_T\_EVENT\_GET\_TIMESTAMP}{(\mbox{event\_instance}, \mbox{event\_timestamp})}
\funcarg{\IN}{event\_instance}{event instance provided to the callback function (handle)}
\funcarg{\OUT}{event\_timestamp}{timestamp the event was observed (integer)}
\end{funcdef}
\mpibindmain{MPI\_T\_event\_get\_timestamp}{(\gb{}MPI\_T\_event\_instance~event\_instance, MPI\_Count~*event\_timestamp)}

\mpifunc{MPI\_T\_EVENT\_GET\_TIMESTAMP} returns the timestamp of when the
event was initially observed by the implementation.  The
\mpiarg{event\_instance} argument identifies the event instance to
query. It is erroneous to provide any other handle to the call than the one
passed by the \MPI/ implementation to the callback function in which the
timestamp is read.

\begin{users}
    An \MPI/ implementation may postpone the call to the user's callback
    function.  In this case, the call to
    \mpifunc{MPI\_T\_EVENT\_GET\_TIMESTAMP} may yield a timestamp in the
    past that is closer to the time the event was initially observed, as
    opposed to a timestamp captured during callback function invocation.
\end{users}

\begin{implementors}
    A high-quality implementation will return a timestamp as close as possible
    to the earliest time the event was observed by the \MPI/ implementation.
\end{implementors}

An event may be raised from different components acting as event sources in
the \MPI/ implementation.  A source in this context is an abstract concept
that helps to define partial ordering of raised events, as each source
provides its own ordering guarantees. A source describes the entity that
raises the event, rather than the origin of the data.

To identify the source of an event instance, the user can query the index
of the source within the corresponding event callback
function invocation.

\begin{implementors}
    An excessive number of event sources may negatively impact performance
    of a tool due to per-source overhead in event handling.
\end{implementors}

\cdeclindex{MPI\_T\_event\_instance}
\begin{funcdef}{MPI\_T\_EVENT\_GET\_SOURCE}{(\mbox{event\_instance}, \mbox{source\_index})}
\funcarg{\IN}{event\_instance}{event instance provided to the callback function (handle)}
\funcarg{\OUT}{source\_index}{index identifying the source (integer)}
\end{funcdef}
\mpibindmain{MPI\_T\_event\_get\_source}{(\gb{}MPI\_T\_event\_instance~event\_instance, int~*source\_index)}

The \mpiarg{event\_instance} argument identifies the event instance to query. It
is erroneous to provide any other event-instance handle to the call than
the one passed by the \MPI/ implementation to the callback function in
which the source is queried.

The \mpiarg{source\_index} argument returns the index of the source of the event instance.
It can be used to query more information on the source using
\mpifunc{MPI\_T\_SOURCE\_GET\_INFO}.

\begin{rationale}
    Event callback function invocations are associated with a source to enable
    chronological processing of events on the tool side, when required,
    while retaining low overhead on the side of the \MPI/ implementation.
\end{rationale}

\subsection{Variable Categorization}
\label{sec:mpit:cat}

\MPI/ implementations can optionally group performance and control
variables into categories to express logical relationships between various
variables.  For example, an \MPI/ implementation could group all control
and performance variables that refer to message transfers in the \MPI/
implementation and thereby distinguish them from variables that refer to
local resources such as memory allocations or other interactions with the
operating system.

Categories can also contain other categories to form a hierarchical
grouping.  Categories can never include themselves, either directly or
transitively within other included categories.  Expanding on the example
above, this allows \MPI/ to refine the grouping of variables referring to
message transfers into variables to control and to monitor message queues,
message matching activities and communication protocols.  Each of these
groups of variables would be represented by a separate category and these
categories would then be listed in a single category representing variables
for message transfers.

The category information may be queried in a fashion similar to the
mechanism for querying variable information.  The \MPI/ implementation
exports a set of $N$ categories via the \MPI/ tool information interface.
If $N=0$, then the \MPI/ implementation does not export any categories,
otherwise the provided categories are indexed from $0$ to $N-1$.  This
index number is used in subsequent calls to functions of the \MPI/ tool
information interface to identify the individual categories.

An \MPI/ implementation is permitted to increase the number of categories
during the execution of an \MPI/ program when new categories become
available through dynamic loading.  However, \MPI/ implementations are not
allowed to change the index of a category or delete it once it has been
added to the set.

Similarly, \MPI/ implementations are allowed to add variables to
categories, but they are not allowed to remove variables from categories or
change the order in which they are returned.

\subsubsection{Category Query Functions}

The following function can be used to query the number of categories, $\mpiarg{num\_cat}$.

\begin{funcdef}{MPI\_T\_CATEGORY\_GET\_NUM}{(\mbox{num\_cat})}
\funcarg{\OUT}{num\_cat}{current number of categories (integer)}
\end{funcdef}
\mpibindmain{MPI\_T\_category\_get\_num}{(\gb{}int~*num\_cat)}

Individual category information can then be queried by calling the
following function:

\begin{funcdef}{MPI\_T\_CATEGORY\_GET\_INFO}{(\mbox{cat\_index}, \mbox{name}, \mbox{name\_len}, \mbox{desc}, \mbox{desc\_len}, \mbox{num\_cvars}, \mbox{num\_pvars}, \mbox{num\_categories})}
\funcarg{\IN}{cat\_index}{index of the category to be queried (integer)}
\funcarg{\OUT}{name}{buffer to return the string containing the name of the category (string)}
\funcarg{\INOUT}{name\_len}{length of the string and/or buffer for \mpiarg{name} (integer)}
\funcarg{\OUT}{desc}{buffer to return the string containing the description of the category (string)}
\funcarg{\INOUT}{desc\_len}{length of the string and/or buffer for \mpiarg{desc} (integer)}
\funcarg{\OUT}{num\_cvars}{number of control variables in the category (integer)}
\funcarg{\OUT}{num\_pvars}{number of performance variables in the category (integer)}
\funcarg{\OUT}{num\_categories}{number of  categories contained in the category (integer)}
\end{funcdef}
\mpibindmain{MPI\_T\_category\_get\_info}{(\gb{}int~cat\_index, char~*name, int~*name\_len, char~*desc, int~*desc\_len, int~*num\_cvars, int~*num\_pvars, int~*num\_categories)}

\textoutdesc{name}{the name of the category}

The routine is required to return a name of at least length one.  This name
must be unique with respect to all other names for categories used by the
\MPI/ implementation.

If any \OUT\ parameter to \mpifunc{MPI\_T\_CATEGORY\_GET\_INFO} is the
\code{NULL} pointer, the implementation will ignore the parameter and not
return a value for the parameter.


\textoutdesc{desc}{the description of the category}

Returning a description is optional.  If an \mpi/ implementation decides
not to return a description, the first character for \mpiarg{desc} must be
set to the null character and \mpiarg{desc\_len} must be set to one at the
return of this call.

The function returns the number of control variables, performance variables
and other categories contained in the queried category in the arguments
\mpiarg{num\_cvars}, \mpiarg{num\_pvars}, and \mpiarg{num\_categories},
respectively.

If the name of a category is equivalent across connected \MPI/ processes, then
the returned description must be equivalent.

\begin{funcdef}{MPI\_T\_CATEGORY\_GET\_NUM\_EVENTS}{(\mbox{cat\_index}, \mbox{num\_events})}
\funcarg{\IN}{cat\_index}{index of the category to be queried (integer)}
\funcarg{\OUT}{num\_events}{number of event types in the category (integer)}
\end{funcdef}
\mpibindmain{MPI\_T\_category\_get\_num\_events}{(\gb{}int~cat\_index, int~*num\_events)}

\mpifunc{MPI\_T\_CATEGORY\_GET\_NUM\_EVENTS} returns the number of event types
contained in the queried category.

\begin{funcdef}{MPI\_T\_CATEGORY\_GET\_INDEX}{(\mbox{name}, \mbox{cat\_index})}
\funcarg{\IN}{name}{the name of the category (string)}
\funcarg{\OUT}{cat\_index}{the index of the category (integer)}
\end{funcdef}
\mpibindmain{MPI\_T\_category\_get\_index}{(\gb{}const~char~*name, int~*cat\_index)}

\mpifunc{MPI\_T\_CATEGORY\_GET\_INDEX} is a function for retrieving the index of a category given a known category name.
The \mpiarg{name} parameter is provided by the caller, and \mpiarg{cat\_index} is returned by the \MPI/ implementation.
The \mpiarg{name} parameter is a string terminated with a null character.

This routine returns \error{MPI\_SUCCESS} on success and returns
\error{MPI\_T\_ERR\_INVALID\_NAME} if \mpiarg{name} does not match the name
of any category provided by the implementation at the time of the call.

\begin{rationale}
This routine is provided to enable fast retrieval of a category index by a
tool, assuming it knows the name of the category for which it is looking.
The number of categories exposed by the implementation can change over
time, so it is not possible for the tool to simply iterate over the list of
categories once at initialization.  Although using \MPI/ implementation
specific category names is not portable across \MPI/ implementations, tool
developers may choose to take this route for lower overhead at runtime
because the tool will not have to iterate over the entire set of categories
to find a specific one.
\end{rationale}

\subsubsection{Category Member Query Functions}
\label{sec:mpit:category_query_functions}

\begin{funcdef}{MPI\_T\_CATEGORY\_GET\_CVARS}{(\mbox{cat\_index}, \mbox{len}, \mbox{indices})}
\funcarg{\IN}{cat\_index}{index of the category to be queried, in the range from $0$ to $\mpiarg{num\_cat}-1$ (integer)}
\funcarg{\IN}{len}{the length of the indices array (integer)}
\funcarg{\OUT}{indices}{an integer array of size \mpiarg{len}, indicating control variable indices (array of integers)}
\end{funcdef}
\mpibindmain{MPI\_T\_category\_get\_cvars}{(\gb{}int~cat\_index, int~len, int~indices[])}

\mpifunc{MPI\_T\_CATEGORY\_GET\_CVARS} can be used to query which control variables are contained in a particular category.
A category contains zero or more control variables.


\begin{funcdef}{MPI\_T\_CATEGORY\_GET\_PVARS}{(\mbox{cat\_index}, \mbox{len}, \mbox{indices})}
\funcarg{\IN}{cat\_index}{index of the category to be queried, in the range from $0$ to $\mpiarg{num\_cat}-1$ (integer)}
\funcarg{\IN}{len}{the length of the indices array (integer)}
\funcarg{\OUT}{indices}{an integer array of size \mpiarg{len}, indicating performance variable indices (array of integers)}
\end{funcdef}
\mpibindmain{MPI\_T\_category\_get\_pvars}{(\gb{}int~cat\_index, int~len, int~indices[])}

\mpifunc{MPI\_T\_CATEGORY\_GET\_PVARS} can be used to query which
performance variables are contained in a particular category.  A category
contains zero or more performance variables.

\begin{funcdef}{MPI\_T\_CATEGORY\_GET\_EVENTS}{(\mbox{cat\_index}, \mbox{len}, \mbox{indices})}
\funcarg{\IN}{cat\_index}{index of the category to be queried, in the range from $0$ to $\mpiarg{num\_cat}-1$ (integer)}
\funcarg{\IN}{len}{the length of the indices array (integer)}
\funcarg{\OUT}{indices}{an integer array of size \mpiarg{len}, indicating event type indices (array of integers)}
\end{funcdef}
\mpibindmain{MPI\_T\_category\_get\_events}{(\gb{}int~cat\_index, int~len, int~indices[])}

\mpifunc{MPI\_T\_CATEGORY\_GET\_EVENTS} can be used to query which
event types are contained in a particular category.  A category
contains zero or more event types.

\begin{funcdef}{MPI\_T\_CATEGORY\_GET\_CATEGORIES}{(\mbox{cat\_index}, \mbox{len}, \mbox{indices})}
\funcarg{\IN}{cat\_index}{index of the category to be queried, in the range from $0$ to $\mpiarg{num\_cat}-1$ (integer)}
\funcarg{\IN}{len}{the length of the indices array (integer)}
\funcarg{\OUT}{indices}{an integer array of size \mpiarg{len}, indicating category indices (array of integers)}
\end{funcdef}
\mpibindmain{MPI\_T\_category\_get\_categories}{(\gb{}int~cat\_index, int~len, int~indices[])}

\mpifunc{MPI\_T\_CATEGORY\_GET\_CATEGORIES} can be used to query which
other categories are contained in a particular category.  A category
contains zero or more other categories.

As mentioned above, \MPI/ implementations can grow the number of categories
as well as the number of variables or other categories within a category.
In order to allow users of the \MPI/ tool information interface to check
quickly whether new categories have been added or new variables or
categories have been added to a category, \MPI/ maintains an update number
that is monotonically increasing during the execution
and is returned by the following function:

\label{sec:mpit:category_changed}
\begin{funcdef}{MPI\_T\_CATEGORY\_CHANGED}{(\mbox{update\_number})}
\funcarg{\OUT}{update\_number}{update number (integer)}
\end{funcdef}
\mpibindmain{MPI\_T\_category\_changed}{(\gb{}int~*update\_number)}

If two calls to this routine return the same update number, it is
guaranteed that the category information has not changed between the two
calls.  If the update number retrieved from the second call is higher, then
some categories have been added or expanded. If the number of changes to
categories exceeds the limit of \mpiarg{update\_number}, an implementation
shall set \mpiarg{update\_number} to the maximum possible value for the type of \mpiarg{update\_number}.

The index values returned in \mpiarg{indices} by
\mpifunc{MPI\_T\_CATEGORY\_GET\_CVARS},
\mpifunc{MPI\_T\_CATEGORY\_GET\_PVARS},
\mpifunc{MPI\_T\_CATEGORY\_GET\_EVENTS}, and
\mpifunc{MPI\_T\_CATEGORY\_GET\_CATEGORIES} can be used as input to
\mpifunc{MPI\_T\_CVAR\_GET\_INFO}, \mpifunc{MPI\_T\_PVAR\_GET\_INFO},
\mpifunc{MPI\_T\_EVENT\_GET\_INFO}, and
\mpifunc{MPI\_T\_CATEGORY\_GET\_INFO}, respectively.

The user is responsible for allocating the arrays passed into the functions
\mpifunc{MPI\_T\_CATEGORY\_GET\_CVARS},
\mpifunc{MPI\_T\_CATEGORY\_GET\_PVARS},
\mpifunc{MPI\_T\_CATEGORY\_GET\_EVENTS}, and
\mpifunc{MPI\_T\_CATEGORY\_GET\_CATEGORIES}.  Starting from array index
$0$, each function writes up to \mpiarg{len} elements into the array.  If
the category contains more than \mpiarg{len} elements, the function returns
an arbitrary subset of size \mpiarg{len}.  Otherwise, the entire set of elements is
returned in the beginning entries of the array, and any remaining array
entries are not modified.

\subsection{Return Codes for the \texorpdfstring{\MPI/}{MPI} Tool Information Interface}
\label{sec:mpit:retcodes}

All procedures defined as part of the \MPI/ tool information interface
return an integer \mpitermni{return code}%
\mpitermindex{return code!tool information interface}%
\mpitermindex{tool information interface!return code} (see Table~\ref{table:tools:mpit:err}) to
indicate whether the function was completed successfully or was aborted.
For the former case, the value \mpiconst{MPI\_SUCCESS} is returned.
In the latter case, the return code indicates the reason for not completing
the routine.  Regardless of whether the return code is \mpiconst{MPI\_SUCCESS} or indicates that the procedure abnormally terminated, the \MPI/ process continues normal execution and does not invoke any \MPI/ error handler.   The \MPI/ implementation is
not required to check all user-provided parameters; if a user passes
invalid parameter values to any routine, the behavior of the implementation
is undefined.

All return codes with the prefix \constskip{MPI\_T\_ERR\_} must be unique values and
cannot overlap with any error codes or error classes returned by the
\MPI/ implementation.  They must also satisfy

$$
0 = \error{MPI\_SUCCESS} < \errorskip{MPI\_T\_ERR\_\XXX/} \leq \error{MPI\_ERR\_LASTCODE}.
$$

\begin{table}[t]
\caption{Return codes used in procedures of the \MPI/ tool information interface}
\label{table:tools:mpit:err} 
\begin{tabularx}{\linewidth}{|l|X|} %ALLOWLATEX%
\hline
\textbf{Return Code} & \textbf{Description} \\
\hline
\hline
\multicolumn{2}{|l|}{{\small Return Codes for All Procedures in the \MPI/ Tool Information Interface}}\\
\hline
\error{MPI\_SUCCESS} & Call completed successfully\\
\errormain{MPI\_T\_ERR\_INVALID} & Invalid or bad parameter value(s)\\
\errormain{MPI\_T\_ERR\_MEMORY} & Out of memory\\
\errormain{MPI\_T\_ERR\_NOT\_INITIALIZED} & Interface not initialized\\
\errormain{MPI\_T\_ERR\_CANNOT\_INIT} & Interface not in the state to be initialized\\ 
\error{MPI\_T\_ERR\_NOT\_ACCESSIBLE} & Requested functionality not accessible\\
\hline
\hline
\multicolumn{2}{|l|}{{\small Return Codes for Datatype Procedures: {\mpiskipfunc{MPI\_T\_ENUM\_\XXX/}}}
\mpifuncindex{MPI\_T\_ENUM\_GET\_INFO}%
\mpifuncindex{MPI\_T\_ENUM\_GET\_ITEM}%
}\\
\hline
\errormain{MPI\_T\_ERR\_INVALID\_INDEX} & The enumeration index is invalid  \\
\hline
\hline
\multicolumn{2}{|l|}{{\small Return Codes for Variable, Category, and Event Query Procedures: {\mpiskipfunc{MPI\_T\_\XXX/\_GET\_\XXX/}}}
\mpifuncindex{MPI\_T\_PVAR\_GET\_INFO}%
\mpifuncindex{MPI\_T\_CVAR\_GET\_INFO}%
\mpifuncindex{MPI\_T\_EVENT\_GET\_INFO}%
\mpifuncindex{MPI\_T\_CATEGORY\_GET\_INFO}%
\mpifuncindex{MPI\_T\_CVAR\_GET\_INDEX}%
\mpifuncindex{MPI\_T\_PVAR\_GET\_INDEX}%
\mpifuncindex{MPI\_T\_CATEGORY\_GET\_INDEX}%
\mpifuncindex{MPI\_T\_EVENT\_GET\_INDEX}%
}\\
\hline
\error{MPI\_T\_ERR\_INVALID\_INDEX} & The variable or category index is invalid\\ 
\errormain{MPI\_T\_ERR\_INVALID\_NAME} & The variable or category name is invalid\\ 
\hline
\hline

\multicolumn{2}{|l|}{{\small Return Codes for Handle Procedures: {\mpiskipfunc{MPI\_T\_\XXX/\_\{ALLOC$|$FREE\}}}}
\mpifuncindex{MPI\_T\_PVAR\_HANDLE\_ALLOC}%
\mpifuncindex{MPI\_T\_CVAR\_HANDLE\_ALLOC}%
\mpifuncindex{MPI\_T\_EVENT\_HANDLE\_ALLOC}%
\mpifuncindex{MPI\_T\_PVAR\_HANDLE\_FREE}%
\mpifuncindex{MPI\_T\_CVAR\_HANDLE\_FREE}%
\mpifuncindex{MPI\_T\_EVENT\_HANDLE\_FREE}%
}\\
\hline
\error{MPI\_T\_ERR\_INVALID\_INDEX} & The variable index is invalid \\ 
\error{MPI\_T\_ERR\_INVALID\_HANDLE} & The handle is invalid\\ 
\errormain{MPI\_T\_ERR\_OUT\_OF\_HANDLES} & No more handles available\\ 
\hline
\hline
%% Reduce the font size so this will fit on one line
\multicolumn{2}{|l|}{{\small Return Codes for Performance Experiment Session Procedures: {\mpiskipfunc{MPI\_T\_PVAR\_SESSION\_\XXX/}}}%
\mpifuncindex{MPI\_T\_PVAR\_SESSION\_CREATE}%
\mpifuncindex{MPI\_T\_PVAR\_SESSION\_FREE}%
}\\
\hline
\errormain{MPI\_T\_ERR\_OUT\_OF\_SESSIONS} & No more sessions available\\ 
\errormain{MPI\_T\_ERR\_INVALID\_SESSION} & Session argument is not a valid session\\
\hline
\hline
\multicolumn{2}{|l|}{{\small Return Codes for Control Variable Access Procedures: {\mpiskipfunc{MPI\_T\_CVAR\_\{READ$|$WRITE\}}}}
\mpifuncindex{MPI\_T\_CVAR\_WRITE}%
\mpifuncindex{MPI\_T\_CVAR\_READ}%
}\\
\hline
\error{MPI\_T\_ERR\_CVAR\_SET\_NOT\_NOW} & Variable cannot be set at this moment\\
\error{MPI\_T\_ERR\_CVAR\_SET\_NEVER} & Variable cannot be set until end of execution\\
\error{MPI\_T\_ERR\_INVALID\_HANDLE} & The handle is invalid\\ 
\hline
\hline
\multicolumn{2}{|l|}{{\small Return Codes for Performance Variable Access and Control Procedures:}}\\

\multicolumn{2}{|l|}{{\small \mpiskipfunc{MPI\_T\_PVAR\_\{START$|$STOP$|$READ$|$WRITE$|$RESET$|$READREST\}}}
\mpifuncindex{MPI\_T\_PVAR\_START}%
\mpifuncindex{MPI\_T\_PVAR\_STOP}%
\mpifuncindex{MPI\_T\_PVAR\_READ}%
\mpifuncindex{MPI\_T\_PVAR\_WRITE}%
\mpifuncindex{MPI\_T\_PVAR\_RESET}%
\mpifuncindex{MPI\_T\_PVAR\_READRESET}%
}\\
\hline
\error{MPI\_T\_ERR\_INVALID\_HANDLE} & The handle is invalid\\ 
\error{MPI\_T\_ERR\_INVALID\_SESSION} & Performance experiment session argument is\flushline invalid\\
\error{MPI\_T\_ERR\_PVAR\_NO\_STARTSTOP} & Variable cannot be started or stopped
                                  (for \mpifunc{MPI\_T\_PVAR\_START} and
                                   \mpifunc{MPI\_T\_PVAR\_STOP})\\
\error{MPI\_T\_ERR\_PVAR\_NO\_WRITE} & Variable cannot be written or reset
                               (for \mpifunc{MPI\_T\_PVAR\_WRITE} and
                               \mpifunc{MPI\_T\_PVAR\_RESET}) \\  
\error{MPI\_T\_ERR\_PVAR\_NO\_ATOMIC} & Variable cannot be read and written\flushline atomically
                              (for \mpifunc{MPI\_T\_PVAR\_READRESET})\\
                                       
\hline
\hline
\multicolumn{2}{|l|}{{\small Return Codes for Source Procedures: {\mpiskipfunc{MPI\_T\_SOURCE\_\XXX/}}}
\mpifuncindex{MPI\_T\_SOURCE\_GET\_INFO}%
\mpifuncindex{MPI\_T\_SOURCE\_GET\_NUM}%
\mpifuncindex{MPI\_T\_SOURCE\_GET\_TIMESTAMP}%
}\\
\hline
\error{MPI\_T\_ERR\_INVALID\_INDEX} & The source index is invalid \\ 
\error{MPI\_T\_ERR\_NOT\_SUPPORTED} & Requested functionality not supported\\
\hline
\hline
\multicolumn{2}{|l|}{{\small Return Codes for Category Procedures: {\mpiskipfunc{MPI\_T\_CATEGORY\_\XXX/}}}
\mpifuncindex{MPI\_T\_CATEGORY\_GET\_CATEGORIES}%
\mpifuncindex{MPI\_T\_CATEGORY\_GET\_CVARS}%
\mpifuncindex{MPI\_T\_CATEGORY\_GET\_PVARS}%
\mpifuncindex{MPI\_T\_CATEGORY\_GET\_INFO}%
}\\
\hline
\error{MPI\_T\_ERR\_INVALID\_INDEX} & The category index is invalid\\
\hline
\end{tabularx}
\end{table} 


\subsection{Profiling Interface}

All requirements for the profiling interface, as described in
Section~\ref{sec:prof}, also apply to the \MPI/ tool information interface.
All rules, guidelines, and recommendations from Section~\ref{sec:prof}
apply equally to procedures defined as part of the \MPI/ tool information
interface.

